% Created 2026-08-01 Sat 21:50
% Intended LaTeX compiler: pdflatex
\documentclass[11pt]{article}

\usepackage[utf8]{inputenc}
\usepackage[T1]{fontenc}
\usepackage{graphicx}
\usepackage{longtable}
\usepackage{wrapfig}
\usepackage{rotating}
\usepackage[normalem]{ulem}
\usepackage{amsmath}
\usepackage{amssymb}
\usepackage{capt-of}
\usepackage{hyperref}
\author{YAMASHITA Takao}
\date{\today}
\title{Hugo blog source}
\hypersetup{
 pdfauthor={YAMASHITA Takao},
 pdftitle={Hugo blog source},
 pdfkeywords={},
 pdfsubject={},
 pdfcreator={},
 pdflang={English}}
\begin{document}

\maketitle
\tableofcontents

\#-\textbf{- mode: org -}-
C-c c b で投稿
C-c C-t でTODOからDONE
C-c C-e H A でOrgからMarkdown変換
hugo server (-D)

\begin{quote}
以下の「ネタ」を、ty07.net（Hugo/ox-hugo, Blowfishテーマ）の
all-posts.org にそのままペーストできる、トップレベル記事エントリの
org-modeサブツリーに変換してください。

\begin{itemize}
\item `* blog` の直下にぶら下がる形式なので、見出しレベルは `**`（レベル2）
\item 見出し直後に以下のPROPERTIESドロワーを付与する（このファイルの既存記事と同じ項目・同じ順序）:
:EXPORT\_HUGO\_BUNDLE: <slug>
:EXPORT\_FILE\_NAME: index
:EXPORT\_DATE: <YYYY-MM-DD HH:MM:SS>
:EXPORT\_HUGO\_TAGS: <タグをスペース区切り>
:EXPORT\_HUGO\_CATEGORIES: <カテゴリ>
:EXPORT\_HUGO\_LASTMOD:
※ Page Bundle化済み（2026-07-30〜）のため、EXPORT\_FILE\_NAMEは常に"index"固定。
記事固有のスラッグはEXPORT\_HUGO\_BUNDLE側に入れること（<slug>はorg-id-new等で生成した
既存記事と同じ形式のUUID/短縮ハッシュを想定）。
※ EXPORT\_HUGO\_LASTMODは空欄のままでよい（\#+hugo\_auto\_set\_lastmod: t のため自動設定される）
\item TODO状態は既存記事の慣習に合わせて付与するか判断する（下書き段階ならTODO、確定ならDONEまたは無し）
\end{itemize}

\begin{itemize}
\item フランクな一人称・だ/である調（このファイルの過去記事のトーンに合わせる）a
\item 新しい感想・解釈・事実を書き足さない。渡された「ネタ」の範囲内で言い回しを整える
\end{itemize}

\begin{itemize}
\item 地の文（本文）は org 記法で統一: 太字 `*string*` / 斜体 `/string/` / コード `=string=` / verbatim `\textasciitilde{}string\textasciitilde{}`
\item `\#+begin\_export hugo` ブロックの中身だけは Markdown 記法で統一: 太字 `**string**` / 斜体 `*string*`
（org記法とMarkdown記法で「アスタリスク1個」の意味が逆転するため、混在させない）
\end{itemize}

\begin{itemize}
\item Perplexity/EllamaなどAI出力の引用は文体を変えず、`\#+begin\_quote` または
Blowfishの `\{\{<details>\}\}` ショートコード（`\#+begin\_export hugo` 経由）で
引用であることを明示する
\item org見出し（`*`〜`****`）は `\#+begin\_quote` 等のブロック内には入れない
\end{itemize}

\begin{itemize}
\item 出力後、`scripts/check\_emphasis.py` での検査を想定した書き方にする
（全角文字に接して \texttt{verbatim} / \texttt{code} / \textbf{bold} が無効化されないよう、
記号の前後にスペースや区切りを意識する。URL・シェルコマンド・カーネルログ・
ショートコード属性内の `` ` `` や `\textasciitilde{}` は対象外）
\end{itemize}

---
ネタ:
<ここに元ネタ・会話内容・ラフな説明文を貼る>
\end{quote}
\section{blog}
\label{sec:org5be21e4}
\subsection{{\bfseries\sffamily DONE} New PC}
\label{sec:org1b2c84b}
息子が iPod nano を買ったところ、今まで使っていたPCでは容量が足りなくなってしまったので、私のPCと交換することにした。交換後に私が使うのは EeePC 901-X という少し古めのUMPCなので、外付けディスクを Linux 用に割り当てて、Live CD 作成マシンとして活用することにした。

Live CD のイメージ作成には、以前は metro を使っていたのだけど、自分しか使わないものなので、もう少し手軽な方法に切り替えることにする。

パーティションレイアウトはこんな感じ。
\subsubsection{fdisk -l}
\label{sec:org5a763e7}
\begin{verbatim}
ディスク /dev/sdc: 160.0 GB, 160041885696 バイト
ヘッド 255, セクタ 63, シリンダ 19457
Units = シリンダ数 of 16065 * 512 = 8225280 バイト
ディスク識別子: 0x1b155f8d

デバイス ブート     始点        終点    ブロック   Id  システム
/dev/sdc1               1        9436    75794638+  83  Linux
/dev/sdc2            9437        9497      489982+  83  Linux
/dev/sdc3            9498       19457    80003700   83  Linux
\end{verbatim}
\subsubsection{/etc/fstab}
\label{sec:org78e1dcb}
\begin{verbatim}
/dev/sde1       /boot       ext3        noauto,noatime  1 2
/dev/mapper/root        /       ext4        noatime     0 1
/dev/mapper/swap    none        swap        sw      0 0
/dev/cdrom      /mnt/cdrom  auto        noauto,ro   0 0
/dev/fd0        /mnt/floppy auto        noauto      0 0
none    /proc/sys/fs/binfmt_misc    binfmt_misc defaults    0 0
/var/tmp/jail/portage   /usr/portage    none        bind,rw     0 0
/var/tmp/jail/local-portage /usr/local/portage  none        bind,rw     0 0
/var/tmp/jail/distfiles /usr/portage/distfiles  none        bind,rw     0 0
\end{verbatim}
\subsubsection{Live CD の起動は Grub で行う}
\label{sec:orgb9dcb18}
\begin{verbatim}
# This is a sample grub.conf for use with Genkernel, per the Gentoo handbook
# http://www.gentoo.org/doc/en/handbook/handbook-x86.xml?part=1&chap=10#doc_chap2
# If you are not using Genkernel and you need help creating this file, you
# should consult the handbook. Alternatively, consult the grub.conf.sample that
# is included with the Grub documentation.

#color white/red black/red
#splashimage=/boot/splash.xpm.gz
splashimage=/boot/grub/bt4.xpm.gz

default 0
timeout 30
password --md5 $1$jBaL5/$pIpowSTX5ip2pDXllzSd90

title=Gentoo Linux (2.6.31-pentoo-r2)
root (hd0,0)
kernel /boot/kernel-genkernel-x86-2.6.31-pentoo-r2 root=/dev/ram0 \
       crypt_root=/dev/sdc3 \
       ramdisk=8192 quiet CONSOLE=/dev/tty1 \
       resume=swap:/dev/mapper/swap init=/linuxrc

initrd /boot/initramfs-genkernel-x86-2.6.31-pentoo-r2

# -- Backtrack4
title BT-4
root (hd0,0)
kernel /boot/bt4/vmlinuz  BOOT=casper boot=casper persistent rw quiet

initrd /boot/bt4/initrd.gz

title=USB stick Pentoo
root (hd0,0)
kernel /boot/kernel-genkernel-x86-2.6.31-pentoo-r2 root=/dev/ram0 \
    root=/dev/ram0 cdroot aufs \
    init=/linuxrc max_loop=256 nokeymap \
    looptype=squashfs loop=/image/root-20091113.squashfs

initrd /boot/initramfs-genkernel-x86-2.6.31-pentoo-r2

#title USB stick Pentoo
#root (hd0,0)
#kernel /boot/pentoo/pentoo \
#    root=/dev/ram0 cdroot aufs changes=/dev/sde2 \
#    init=/linuxrc max_loop=256 nokeymap \
#    looptype=squashfs loop=/pentoo/image-2009.squashfs
#initrd /boot/pentoo/pentoo.igz

#title USB stick Pentoo
#root (hd0,0)
#kernel /boot/kernel-genkernel-x86-2.6.29-pentoo-r2 \
#    root=/dev/ram0 cdroot aufs changes=/dev/sdc2 \
#    init=/linuxrc max_loop=256 nokeymap \
#    looptype=squashfs loop=/pentoo/image.squashfs
#
#initrd /boot/initramfs-genkernel-x86-2.6.29-pentoo-r2

title grub-install
lock
install (hd0,0)/boot/grub/stage1 d (hd0) (hd0,0)/boot/grub/stage2 p (hd0,0)/boot/grub/grub.conf

title Other Operating System - Microsoft Windows XP
lock
    rootnoverify (hd0,0)
    makeactive
    chainloader +1

# vim:ft=conf:
\end{verbatim}
\subsubsection{Live CD の作成は chroot 環境で行う}
\label{sec:org80e3040}
\begin{itemize}
\item \href{http://github.com/ac1965/config-ac1965/blob/master/etc/skel/script/in.sh}{CHROOT環境設定}
\item \href{http://github.com/ac1965/config-ac1965/blob/master/etc/skel/script/in.sh}{コンパイル用}
\item \href{http://github.com/ac1965/config-ac1965/blob/master/etc/skel/script/chroot.sh}{CHROOT}
\item \href{http://github.com/ac1965/config-ac1965/blob/master/etc/skel/script/mkrootimg.sh}{MAKE ROOT-IMAGE}
\end{itemize}

こんな感じで使っている。CHROOT は私の環境では "/var/tmp/jail/squashfs-root" にしている。
\begin{verbatim}
# ./in.sh
# cp chroot.sh ${CHROOT}/tmp
# chroot ${CHROOT} /tmp/chroot.sh
# ./out.sh
# ./mkrootimg.sh
\end{verbatim}
\subsection{{\bfseries\sffamily DONE} awesome window manager}
\label{sec:orgc4abc8b}
\href{http://awesome.naquadah.org/}{awesome} の設定を見直してみた。見た目にはあまり変わっていないように見えるかもしれないけど、tag（他のウィンドウマネージャでいうところのワークスペース）を動的に生成する shifty というのを使ってみている。

\begin{itemize}
\item git://github.com/bioe007/awesome-configs.git
\end{itemize}
\subsection{{\bfseries\sffamily DONE} Guruplug Std.}
\label{sec:orge6f97ca}
\subsubsection{今なら(2/28/2010まで)、JTAGモジュールが無料の組み込み}
\label{sec:org57b679c}
今なら（2/28/2010まで）、JTAGモジュールが無料でついてくる組み込みLinuxキットの \href{http://www.globalscaletechnologies.com/p-31-guruplug-server-standard.aspx}{guruplug std.} を注文した。4月から配送が始まるらしいので、5月までには届いてくれるだろうか。

\begin{itemize}
\item \url{http://www.plugcomputer.org/}
\item \url{http://plugcomputer.org/plugwiki/index.php?title=Main\_Page}
\item \url{http://www.openplug.org/data/}
\end{itemize}
\subsection{{\bfseries\sffamily DONE} xmonad window manager}
\label{sec:org4922057}
awesome が最高だと思っていたけど、\href{http://xmonad.org/}{Xmonad} はさらに良かった。キーボードだけで操作できるし、awesome と違って設定ファイルが多少おかしくなっても以前の設定のまま動いてくれるので、使う側としては安心感がある。awesome で気に入っていたものは、だいたい xmonad 側にも揃っている。bioe007 が書いた awesome の revelation がお気に入りだったので、同じようなものが xmonad にもないか、いま探しているところ。

\begin{itemize}
\item git clone git://github.com/jelly/dotfiles.git
\end{itemize}
\subsection{{\bfseries\sffamily DONE} Xmonad for Backtrack4}
\label{sec:org20a848b}
backtrack4 に xmonad を入れてみた。

\href{http://www.lares.dti.ne.jp/\~jargon/uploads/backtrack4-xmonad.png}{{[}[http://www.lares.dti.ne.jp/\textasciitilde{}jargon/uploads/backtrack4-xmonad.png}]]

\begin{itemize}
\item \$HOME/.xmonad/xmonad.hs
\end{itemize}

\begin{verbatim}
import XMonad
import XMonad.Config.Kde
import qualified XMonad.StackSet as W -- to shift and float windows

main = xmonad $ kdeConfig

 { modMask = mod4Mask -- use the Windows button as mod
 , manageHook = manageHook kdeConfig <+> myManageHook
 }
 where
   myManageHook = composeAll . concat $
     [ [ className   =? c --> doFloat           | c <- myFloats]
     , [ title       =? t --> doFloat           | t <- myOtherFloats]
     , [ className   =? c --> doF (W.shift "2") | c <- webApps]
     , [ className   =? c --> doF (W.shift "3") | c <- ircApps]
     ]
   myFloats      = ["MPlayer", "Gimp"]
   myOtherFloats = ["alsamixer"]
   webApps       = ["Firefox-bin", "Opera"] -- open on desktop 2
   ircApps       = ["Ksirc"]                -- open on desktop 3
\end{verbatim}

\begin{itemize}
\item 設定
\end{itemize}

\begin{verbatim}
sudo apt-get install language-pack-ja
sudo apt-get install scim scim-anthy kasumi
sudo apt-get install ttf-vlgothic
sudo apt-get install kde-i18n-ja
sudo apt-get install xmonad

cat $HOME/.kde3/env/set_window_manager.sh
KDEWM=/usr/bin/xmonad
chmod +x $HOME/.kde3/env/set_window_manager.sh
\end{verbatim}
\subsection{{\bfseries\sffamily DONE} USB Encrypt for Backtrack}
\label{sec:org18fe79c}
\href{http://www.backtrack-linux.org/forums/}{backtrack4のフォーラム}を読んでいたら、質問だけして自分では調べようとしない人を見かけたのでフォローしておいた。内容は distro の iso を更新する方法についてだったのだけど、\href{http://www.offensive-security.com/blog/backtrack/customising-backtrack-live-cd-the-easy-way/}{bt4-customise.sh}を使えばできるとのこと。

さらに手を加えたものをここに置いておく。squashfs を LUKS で暗号化するものになっている。
\subsection{{\bfseries\sffamily DONE} Guruplug 届くまでにやったこと}
\label{sec:org2349de2}
\subsubsection{guruplug が届く前に環境を作っておこう}
\label{sec:org562ad63}
\subsubsection{1日目}
\label{sec:org42f26f9}
クロスコンパイル環境やルートファイルシステムまわりの作業メモ。\href{http://www.plugcomputer.org/plugwiki/index.php/GuruPlug\%22}{GuruPlug Wiki} に自分なりの補足を加えたものも含んでいる。
\begin{enumerate}
\item 参考にしたポイント
\label{sec:org939e62a}
\begin{itemize}
\item \href{http://www.gentoo.org/proj/en/base/embedded/cross-development.xml}{cross-development}
\item \href{http://computingplugs.com/index.php/Building\_a\_custom\_kernel}{Building
Kernel}
\item \href{http://www.plugcomputer.org/plugwiki/index.php/GuruPlug}{GuruPlug}
\end{itemize}

クロスコンパイルは \href{http://www.gnu.org/software/hurd/}{GNU Hurd} で遊んで以来ずっと触っていなかったけれど、Gentoo だとかなり楽にできる。
\item クロスコンパイル環境のセットアップ
\label{sec:orge3e42d4}
\begin{verbatim}
sudo emerge -av portage-utils crossdev
sudo crossdev armv5tel-softfloat-linux-gnueabi
\end{verbatim}
\item U-Boot とカーネルセットアップ
\label{sec:org1762c8f}
ここは Wiki に書いてあるとおりに進めればOK。最後に mkimage を PATH の通っているところにコピーしておくのを忘れずに。これをやらないとカーネル（uImage）のビルドでこけてしまう。

\begin{verbatim}
sudo git clone git://git.denx.de/u-boot-marvell.git u-boot-marvell.git
sudo cd u-boot-marvell.git
sudo git checkout -b testing origin/testing
sudo make mrproper
sudo make guruplug_config
sudo make u-boot.kwb CROSS_COMPILE=armv5tel-softfloat-linux-gnueabi-
sudo cp tools/mkimage /usr/bin
sudo wget http://www.plugcomputer.org/plugwiki/images/8/81/Guruplug-patchset.tar.bz2
sudo tar xfj Guruplug-patchset.tar.bz2
sudo git clone git://git.kernel.org/pub/scm/linux/kernel/git/stable/linux-2.6.33.y.git
sudo for p in guruplug-patchset/*; do patch -p1 -E -d linux-2.6.33.y.git < $p; done
sudo cd linux-2.6.33.y.git
sudo make CROSS_COMPILE=armv5tel-softfloat-linux-gnueabi- ARCH=arm clean
sudo make CROSS_COMPILE=armv5tel-softfloat-linux-gnueabi- ARCH=arm guruplug_defconfig
sudo make CROSS_COMPILE=armv5tel-softfloat-linux-gnueabi- ARCH=arm uImage
sudo make CROSS_COMPILE=armv5tel-softfloat-linux-gnueabi- ARCH=arm modules
sudo make CROSS_COMPILE=armv5tel-softfloat-linux-gnueabi- ARCH=arm modules_install INSTALL_MOD_PATH=/usr/armv5tel-softfloat-linux-gnueabi
\end{verbatim}
\item mtd-toolsのセットアップ
\label{sec:orgb71d386}
git リポジトリ（git://git.infradead.org/mtd-utils.git）を使ってみることにする。ubifs でルートファイルシステムを構築するための準備。

\begin{verbatim}
sudo ebuild /usr/portage/sys-fs/mtd-utils/mtd-utils-99999999.ebuild compile
sudo ebuild /usr/portage/sys-fs/mtd-utils/mtd-utils-99999999.ebuild merge
sudo ebuild /usr/portage/sys-fs/mtd-utils/mtd-utils-99999999.ebuild clean
\end{verbatim}
\item ルートファイルシステムのセットアップ
\label{sec:org46d3fbe}
crossdev で作ったクロスコンパイラで基本システム（stage1）を構築しようとしたところ、`cpio' のコンパイルで失敗してしまった。

今日は眠いのでここまでにしておく。

\begin{verbatim}
cd /usr/armv5tel-softfloat-linux-gnueabi
mkdir etc
cd etc
ln -s /usr/portage/profiles/default/linux/arm/10.0 make.profile
cp /tmp/make.conf-arm make.conf
armv5tel-softfloat-linux-gnueabi-emerge -av @system
\end{verbatim}

\begin{itemize}
\item 参考：/usr/armv5tel-softfloat-linux-gnueabi/etc/make.conf
\end{itemize}

あくまで自分の環境用の設定なので、同じ内容にする必要はない。ROOT="\ldots{}" のデフォルトは "/" なので、きちんと指定しておかないと上書きされてしまうので注意。あと、クロスコンパイルの場合は CBUILD の指定も必要になる。
\end{enumerate}
\subsubsection{2日目}
\label{sec:orgc858d9b}
\begin{enumerate}
\item cpioのコンパイルエラーは stat(2) の引数が指定されていなかっただけだった
\label{sec:org1e68f7e}
野良ebuildを挟んで解決。

\begin{itemize}
\item src/filtypes.h の patch
\end{itemize}

\begin{verbatim}
diff -urN cpio-2.11-/src/filetypes.h cpio-2.11/src/filetypes.h
--- cpio-2.11-/src/filetypes.h  2010-04-15 18:03:16.639694128 +0900
+++ cpio-2.11/src/filetypes.h   2010-04-15 18:04:03.545691921 +0900
@@ -82,4 +82,4 @@
 #define lstat stat
 #endif
 int lstat ();
-int stat ();
+int stat (const char *path, struct stat *buf);
\end{verbatim}
\item 野良ebuild cpio-2.11-r1.ebuild
\label{sec:org71177a4}
\end{enumerate}
\subsubsection{3日目}
\label{sec:org50ce199}
\begin{enumerate}
\item クロス環境での stage1 がエラーで出来ていない状態
\label{sec:org10ca299}
肝心のguruplug本体がまだ届いていないので、焦らずゆっくり進めることにする。あとで distcc を使ってホストマシン側でコンパイルできるように、クロスコンパイル環境を作り直しておこう。

一度クロス環境を消して、ホスト環境とバージョンを合わせておくことにする。

SYSROOT を検索する必要があるため、binutils は 2.19.51.0.12 より新しいものにする。参考：\url{http://bugs.gentoo.org/275666}

\begin{verbatim}
sudo crossdev --clean armv5tel-softfloat-linux-gnueabi
sudo emerge -p binutils linux-headers gcc glibc
....
[ebuild   R   ] sys-devel/binutils-2.20.1
[ebuild   R   ] sys-kernel/linux-headers-2.6.29
[ebuild   R   ] sys-devel/gcc-4.3.3-r2
[ebuild   R   ] sys-libs/glibc-2.9_p20081201-r2
...
sudo crossdev -t armv5tel-softfloat-linux-gnueabi \
    --b 2.20.1 \
    --k 2.6.29 \
    --g 4.3.3-r2 \
    --l 2.9_p20081201-r2 \
    --ex-gdb
\end{verbatim}
\end{enumerate}
\subsubsection{4日目}
\label{sec:orgefd1e5b}
\begin{enumerate}
\item クロスコンパイル環境で ARM(armv5tel)用の gentoo stage1
\label{sec:org8850d63}
ができたので動作確認をした
:CUSTOM\_ID: クロスコンパイル環境で-armarmv5tel用の-gentoo-stage1-ができたので動作確認をした

perl と Linux-PAM についてはクロスコンパイルがうまくいかなかったので、こちらは別途自分で用意することにする。
\item \url{http://www.gentoo.org/proj/en/base/embedded/handbook/?part=1\&chap=5}
\label{sec:org4e18bce}
\end{enumerate}
\subsection{{\bfseries\sffamily DONE} Xorg}
\label{sec:orgd61214c}
GuruPlug がなかなか届かないので、待っている間に Funtoo と xorg-1.8 まわりを整理してみた。
\subsubsection{boot-update}
\label{sec:org97d8672}
grubの設定を助けてくれるツールのようだ。\href{http://www.funtoo.org/en/funtoo/core/boot/}{boot-update} というもので、grub-1.97+を試してみたけれど、multiboot にはまだきちんと対応していないみたい。/etc/boot.conf に次のように書いておいて `boot-update` と打つだけで済むので楽ではあるものの、"Backtrack 4" のほうは利用側の root-fs が無理やり埋め込まれてしまっているので、少し手直しが必要になる。/boot/grub/grub.cfg を手で修正すれば大丈夫。

\begin{verbatim}
boot {
    path /boot
    generate grub
    default "Funtoo Linux"
    timeout 3
}

"Funtoo Linux" {
    kernel /kernel-genkernel-x86[-v]
    initrd initramfs-genkernel-x86[-v]
    params += crypt_root=/dev/sdc2 root_keydev=/dev/sde1 root_key=/keyfile
    params += dolvm real_root=/dev/mapper/LVG-root
    params += i915.modeset=1 fbcon=map:1
    params += ramdisk=8192 quiet init=/linuxrc
}

"Backtrack 4" {
    kernel /bt4/vmlinuz
    initrd /bt4/initrd.gz
    params += BOOT=casper boot=casper persistent rw quiet
    params += real_root=auto
}
\end{verbatim}
\subsubsection{xorg-1.8}
\label{sec:org14014b7}
しばらく放置していたので、そろそろと思い試してみることにした。
\begin{enumerate}
\item MASKを外す
\label{sec:orgd1b496a}
\begin{verbatim}
echo 'x11-base/xorg-server' >> /etc/portage/package.unmask
echo 'x11-base/xorg-server * ~* **' >> /etc/portage/package.keywords/x11-base
\end{verbatim}
\end{enumerate}
\subsubsection{emege xorg-server}
\label{sec:org7c40121}
USE="udev -hal" で emerge してみたところ、キーボードとマウスが認識されなくなってしまった。ある程度予想はしていたので、sshd を立ち上げておいて、別端末から pkill して対処した。

参考：\href{http://body0r.wordpress.com/2010/04/16/xorg-udev-toggle/}{こちらの記事}
\subsubsection{emerge udev}
\label{sec:orga6c7362}
MASKを外して、udev をアップデートする。

\begin{verbatim}
echo 'sys-fs/udev' >> /etc/portage/package.unmask
echo 'sys-fs/udev * ~* **' >> /etc/portage/package.keywords/sys-fs
emerge -u udev
\end{verbatim}
\subsubsection{udevルールの追加}
\label{sec:org02db98f}
/usr/share/X11/xorg.conf.d がシステム設定になっているので、/etc/X11/xorg.conf.d 側にキーボードとマウスの設定を追加しておく。

これで無事解決。

\begin{verbatim}
# cat /etc/X11/xorg.conf.d/10-keyboard.conf
Section "InputClass"
        Identifier "Keyboard"
        Driver "evdev"
        MatchIsKeyboard "on"
        Option "xkbmodel" "jp106"
        Option "xkblayout" "jp"
EndSection

# cat /etc/X11/xorg.conf.d/20-synaptics.conf
Section "InputClass"
    Identifier "Touchpad"
    Driver "synaptics"
    MatchIsTouchpad "on"
    Option "SHMConfig" "true"
    Option "MinSpeed" "0.20"
    Option "MaxSpeed" "0.60"
    Option "AccelFactor" "0.020"
    Option "HorizEdgeScroll" "true"
    Option "HorizScrollDelta" "100"
    Option "VertEdgeScroll" "true"
    Option "VertScrollDelta" "100"
    Option "TapButton1" "1"
EndSection
\end{verbatim}
\subsection{{\bfseries\sffamily DONE} Guruplug Arrived}
\label{sec:org2c0d400}
出遅れてしまった感は否めないけれど（まぁ人柱になってくれた人がたくさんいたので、その点は心強かった）、やっと Guruplug Std. が届いた。

もしかすると Plus の注文と間違われていたのかもしれない。一昨日 GlobalScale 社に「まだ発送されないんですか？」というメールを送った際、「自分が頼んだのは Plus じゃないですよ」と付け加えたところ、翌日には「明日発送します」という返信が届いた。

熱くなりやすいのは Plus と同じかもしれない。どこまで熱くなるかはまた次回試してみようと思う（週末にでも。でも、何かイベントがあったような気もするけど……）。
\subsubsection{おきまりの first impression ということで、お写真を}
\label{sec:org670a2bc}
以前のブログサイトにアップロードしていた画像なので、こちらには残っていない（w）
\subsubsection[JTAG ボード を繋ぎ、FTDI USB シリアルデバイス(例ではttyUSB0)を確認したら]{JTAG ボード \footnote{Emacsのエディタ機能は大量のLispで作られているため処理速度は遅く重たい。ネイティブコードとは実行する都度読んで解釈していたLispコードをバイナリコードにバックグラウンドで自動的にコンパイルする機能で、次回から最速で動作することが期待される。\label{org3132e76}} を繋ぎ、FTDI USB シリアルデバイス(例ではttyUSB0)を確認したら}
\label{sec:orgb38f801}
\begin{verbatim}
usb 2-2: FTDI USB Serial Device converter now attached to ttyUSB0
\end{verbatim}

\texttt{screen /dev/ttyUSB0 115200} として `screen` コマンドでターミナルを立ち上げておく。
\subsection{{\bfseries\sffamily DONE} Guruplug Subquent}
\label{sec:org0671de9}
\subsubsection{guruplug が届いてから遊んだこと}
\label{sec:org90ccdc3}
\begin{itemize}
\item U-Bootのアップデート
\item Kernelのアップデート
\item Debian RFS から Gentoo RFS に変更
\end{itemize}

この3つはすでに確認済み。まぁ毎回やることでもないので、いったんこのままにしておく。

wiki に書いてあることで十分だった。Gentoo RFS については、stage3 を持ってくるのが一番手軽かなと思う。インストール自体は x86 のときと同じ感覚でできる。

参考にしたポイント

\begin{itemize}
\item \url{http://mark.ossdl.de/2009/09/gentoo-on-the-sheevaplug/}
\end{itemize}
\subsubsection{進行形}
\label{sec:org8f34b63}
\begin{itemize}
\item Pyblosxom + Paste + Nginx
\end{itemize}

python アプリ関係のソフトに "arm" や "\textasciitilde{}arm" の KEYWORD がついていなくて、ちょっと困ってしまった。

KEYWORD のチェックを無効にしてから Nginx をインストールし、動作するところまでは確認できた。ただ、評判どおりの（軽い/速い）動作にはなっていないので、いまも調整中。

Nginx + paster + Pyblosxom の組み合わせでは動作を確認できたのだけど、wbgpager だけはうまく動かなかった。

原因を追いかけるのが面倒だったので、いまのところは Apache + mod\_wsgi + Pyblosxom で動かしている。\textit{{[}2010-07-19 Mon]}
\subsubsection{将来}
\label{sec:orgcd44a49}
\begin{itemize}
\item HoneyPod の実装
\end{itemize}

GuruPlug を使ったアイデアとして、やりたいのは HoneyPod の実装だ。体感的にはそんなに遅いマシンでもないので期待しているのだけど、Nginx がうまく動作していない分、まだ実現には遠い。

あとは、カーネルが変わっても iptables がきちんと動作してくれるかどうか。

iptables は問題なく動作した。honeyd も動いている。あとはルータ側で honeyd の virtual-honeypots に流し込むだけ。

ただ、mwcollect が SIGSEGV で落ちてしまっている。\textit{{[}2010-07-19 Mon]}
\subsubsection{最後に前回の宿題}
\label{sec:org46e3d68}
耐久テストとしては、届いてから（今週月曜日）カーネルを変更したりしていたので実質4日間くらいの稼働になるけれど、いまのところ全く問題ない。熱暴走はまだ起きていないものの、USBスティックのほうが熱くなっている感じがする。

本体の発熱については、某掲示板で騒がれているほどひどくはない。話題になっているのは GuruPlug PLUS（ethx2 + eSTA のほう）だけのようだ。
\subsection{{\bfseries\sffamily DONE} 四季・ユートピアノ 中尾幸世さん}
\label{sec:org8f79616}
「四季・ユートピアノ」とか「中尾幸世」っていう女優（朗読家）、知ってる人はどれくらいいるだろう。

先週の金曜日（2010-07-30）、45歳の誕生日を自分なりに祝おうと思って、渋谷のユーロスペースでやっていた上映会「佐々木昭一郎というジャンル」の最終日に足を運んだ。中尾幸世さんのトークもあるらしいと聞いていたので、前売券は早めに確保しておいた。おかげで満員のユーロスペースでもいい席に座れて、「四季・ユートピアノ」を懐かしい気持ちで観ることができた。

\begin{center}
\includegraphics[width=.9\linewidth]{20100730266.jpg}
\end{center}

\href{https://github.com/yt-dlp/yt-dlp}{yt-dlp} みたいなダウンローダーで映像を落とすのは、たしか違法だったはず🤔　アップロードしたり売ったりするのが違法だというのは覚えているんだけど、細かいところはあやふやになってしまった。
\subsection{{\bfseries\sffamily DONE} Hexo でハロー}
\label{sec:orgbb3f53c}
Welcome to \href{http://hexo.io}{Hexo}!
\subsubsection{環境}
\label{sec:org1215e50}
\begin{itemize}
\item OS X Yosemite 10.10.3
\item Node.js
\end{itemize}

\begin{verbatim}
git clone git://github.com/creationix/nvm.git ~/.nvm
test -f ~/.nvm/nvm.sh && source ~/.nvm/nvm.sh
nvm ls-remote
nvm install 0.12.2
node -v
v0.12.2
\end{verbatim}
\subsubsection{前準備}
\label{sec:org6a79fb3}
github にリポジトリを作っておく。中身は空っぽで構わない。

例：Mac の brew で hub をインストールしている場合

\begin{verbatim}
cd ~/Documents/github
mkdir ac1965.github.io
cd ac1965.github.io
hub init
hub create ac1965/ac1965.github.io
\end{verbatim}
\subsubsection{Hexo インストール}
\label{sec:org99bf0bd}
\begin{verbatim}
$ npm install -g hexo
\end{verbatim}
\begin{enumerate}
\item 準備（ブログ用のディレクトリを作る）
\label{sec:org1a6122d}
hexo init フォルダ名、という形で使う。

フォルダを指定しないとカレントディレクトリに作られてしまうので、きちんと指定しておこう。

\begin{verbatim}
cd ~/Documents
$ hexo init site
$ cd site
$ npm install
\end{verbatim}
\item 動作確認
\label{sec:org136ecf7}
\begin{verbatim}
$ hexo server
\end{verbatim}

この状態でウェブブラウザから\href{http://localhost:4000/}{http://localhost:4000}にアクセスして確認してみる。
\item Github にデプロイしてみる
\label{sec:org344afe2}
Hexo の 3.0 以降は、ぐぐって出てきた情報とはデプロイの仕組みが少し変わっているようだ。

\begin{verbatim}
$ hexo deploy
ERROR Deployer not found: github
\end{verbatim}

要するにデプロイヤが github から git に変わっていて、追加でモジュールをインストールする必要があった。

\begin{verbatim}
$ npm install hexo-deployer-git --save
\end{verbatim}

詳しくは：\url{https://github.com/hexojs/hexo/issues/1040}

参考：\_config.ymlの deploy部分

\begin{verbatim}
deploy:
  type: git
  repo: git@github.com:ac1965/ac1965.github.io.git
  branch: master
\end{verbatim}
\item テーマを変えてみよう
\label{sec:org5b2c603}
\href{http://hexo.io/themes/}{テーマ一覧}から気に入ったものを探してみるのもいい。

\begin{verbatim}
$ git clone https://github.com/joyceim/hexo-theme-apollo.git themes/apollo
\end{verbatim}

参考：\_config.ymlの languageとthemes 部分

\begin{verbatim}
language: defualt
...
theme: apollo
\end{verbatim}
\item ポスト
\label{sec:org9e0c7cb}
\begin{verbatim}
$ hexo new "My New Post"
INFO  Created: ~/Documents/site/source/_posts/My-New-Post.md
\end{verbatim}

source/\_posts/My-New-Post.md をエディタで編集する（\href{https://help.github.com/articles/github-flavored-markdown/}{GFM}：Github Flavored Markdown 形式）。

詳しくは： \href{http://hexo.io/docs/writing.html}{Writing}
\item 静的ファイルの作成
\label{sec:orgd2a7356}
\begin{verbatim}
$ hexo generate
$ hexo generate --deploy
$ hexo deploy --generate
\end{verbatim}

詳しくは： \href{http://hexo.io/docs/generating.html}{Generating}
\end{enumerate}
\subsection{{\bfseries\sffamily DONE} Using Packer for Archassult}
\label{sec:org90dd33e}
\subsubsection{すこしまとめてみる}
\label{sec:org9001b3c}
\begin{enumerate}
\item 前提
\label{sec:org7194231}
\begin{itemize}
\item base は ArchLinux ISO からつくる
\item ArchAssault の環境は repository 掘るところ位まで
\item もちろん
\href{https://www.virtualbox.org}{virtualbox}、\href{https://www.vagrantup.com}{vagrant}、\href{https://www.vagrantup.com}{packer}
の環境は作成済み
\end{itemize}
\end{enumerate}
\subsubsection{じつのところ}
\label{sec:org4080e55}
Packer の細かい設定はまだよくわかっていない。Vagrant の box を作るためのパッケージだということは知っているけれど。というわけで、ベースの box さえ作ってしまえば、あとは \texttt{vagrant up} でどんどん環境をデプロイしていくほうが良さそうだと思い始めている。
\subsubsection{作成}
\label{sec:org91bba29}
\begin{verbatim}
$ git clone https://github.com/ac1965/archassault-packer.git
$ cd archassult-packer
$ packer build -only=virtualbox-iso template.json
$ vagrant box add archassault packer_archassault_virtualbox.box
$ vagrant init archassult
$ vagrant up
\end{verbatim}
\begin{enumerate}
\item 手をつけているところ
\label{sec:org9b6f5a6}
Cuckoo環境の半自動生成
\begin{itemize}
\item \href{http://www.cuckoosandbox.org}{Cuckoo}で簡易マルウェア自動解析をしてみたい
\item \href{http://cuckoo.readthedocs.org/en/latest/installation/guest/}{ゲストOSの準備}としては、あらかじめ \href{https://www.modern.ie/ja-jp}{ModernIE} から Linux 版の Virtualbox イメージをダウンロードしておき、\href{https://www.python.org/getit/}{Python} と \href{http://www.pythonware.com/products/pil/}{PIL(Python Image Library)} をインストールしておく。cuckoo の agent.py は起動時に自動実行できるよう agent.pyw にリネームしておく。
\end{itemize}

\begin{verbatim}
reg.exe ADD "HKEY_LOCAL_MACHINE\Software\Microsoft\Windows\CurrentVersion\Run" -v "Agent" -t REG_SZ -d "C:\Python27\agent.pyw"
\end{verbatim}
\end{enumerate}
\subsubsection{参考}
\label{sec:org0f5551b}
\begin{itemize}
\item \href{https://www.archlinux.org}{Archlinux}
\item \href{https://www.archassault.org}{ArchAssault}
\item \href{https://github.com/elasticdog/packer-arch}{Archlinux用のPackerテンプレ}
\item \href{https://github.com/ac1965/archassault-packer}{オレオレ Packer}
\end{itemize}
\subsection{{\bfseries\sffamily DONE} 健康診断}
\label{sec:org3260a45}
\subsubsection{定期検診}
\label{sec:org90a5bea}
今日は60日ぶりの定期検診。

薬もいただかないといけない。ジェネリックばかりだけど。\begin{center}
\includegraphics[width=.9\linewidth]{medicine.jpg}
\end{center}

血圧は正常だった。主治医からは「血圧だけは正常ですね」なんて言われる始末（笑）\begin{center}
\includegraphics[width=.9\linewidth]{blood-pressure-result.jpg}
\end{center}
\subsection{{\bfseries\sffamily DONE} What Blog?}
\label{sec:orgacec0ad}
\subsubsection{ブログってw}
\label{sec:orge6ce88a}
毎日書き続けるものだから、筆不精な自分には向いていないというのはよくわかっている。紙の日記も三日坊主だったし。

そもそも、ブログというのは作っては壊しての繰り返しで、結局ログとしてまともに残っていないことのほうが多いものだ。

自分が日記をつけるとしたら、きっとのび太の日記みたいになるんだろうな。

「眠い」「お腹が空いた」

とか、そんなことばかり。でもよく考えると、のび太が今の時代にいたら、Twitterのヘビーユーザーになっていたかもしれない（笑）
\subsection{{\bfseries\sffamily DONE} Hugo}
\label{sec:orga0adc14}
\subsubsection{Hexo から Hugo}
\label{sec:org31152f3}
サイトジェネレータを \href{https://gohugo.io}{Hugo} に変えてみた。インストールといっても homebrew を使うので簡単 :grin:

\href{https://www.webfx.com/tools/emoji-cheat-sheet/}{EMOJI-CHEAT-SHEET}
\subsubsection{インストール}
\label{sec:org6e8c2a2}
\href{https://brew.sh/index\_ja}{Homebrew} はインストール済みという前提で進める。
\begin{verbatim}
$ brew update
$ brew install hugo
\end{verbatim}
テーマは \href{https://conao3.com}{conao3さん} の \href{https://github.com/conao3/anatole-ext}{anatole-ext} にすることにした。
\begin{verbatim}
$ hugo new site hoge
$ cd hoge
$ git clone https://github.com/conao3/anatole-ext themes/anatole-ext
$ hugo new posts/first-article.md
\end{verbatim}
これで hoge ディレクトリーが作成され、その配下に雛形が作成される。
\subsubsection{自分用のメモ}
\label{sec:org9c7b4b8}
\begin{enumerate}
\item 展開
\label{sec:org03c6d08}
\begin{verbatim}
$ git clone git@github.com:ac1965/hugo.git ./hugo-blog
$ cd hugo-blog
$ git clone git@github.com:ac1965/anatole-ext themes/anatole-ext
$ hugo
\end{verbatim}
hugo コマンドを実行すると、public ディレクトリの下にコンテンツが生成される。
\item `deploy.sh'
\label{sec:org303251a}
hexo と大きく違うのは、デプロイの仕組みが利用者まかせになっているところ。

私の環境では本体用と公開用でリポジトリを分けていて、本体は hugo-blog、公開用は ac1965.github.com として管理している。
\end{enumerate}
\subsection{{\bfseries\sffamily DONE} iPhone 12 Pro}
\label{sec:org6ad6552}
機種変更のタイミングということで、iPhone を 8plus から 12pro に変えた。長らく8plusを使い倒してきたので、正直このタイミングを逃したらズルズル使い続けていた気がする。

8plus からの変化で大きいのは、なんといっても顔認証になったことだ。Touch ID に体が慣れきっていたせいで、最初の数日はついホームボタンを探して指を当てようとしてしまい、そのたびに「あ、もう無いんだった」と思い出す羽目になった。それでも数日も使えばすっかり慣れて、今ではロック解除のたびに画面を見るだけで済むのがむしろ自然になっている。

それより地味に嬉しいのは、ひとまわり小さくなって持ちやすくなった点で、これがなかなかいい感じだ。8plus はどうしても片手操作がしんどい場面が多かったんだけど、12pro はポケットへの収まりも手のひらへのフィット感も良くなって、日常のちょっとした場面でのストレスが減った気がする。

アプリと Wallet はひととおり復元できたけど、一番面倒だったのは ○○Pay 系のアプリで、これは再認証が必要なものがいくつかあった。バックアップから機種変更すればそのまま使えるだろうと高を括っていたら、決済系だけは律儀に本人確認を求めてくるあたり、セキュリティ的には妥当なんだろうけど地味に手間だった。
\subsubsection{箱から開けた}
\label{sec:org47d34d0}

\begin{center}
\includegraphics[width=.9\linewidth]{1st-shot.png}
\end{center}

開封の瞬間。何度やってもこの瞬間だけはちょっとテンションが上がる。
\subsubsection{PCだけの光}
\label{sec:orgfa6ba09}

\begin{center}
\includegraphics[width=.9\linewidth]{2nd-shot.png}
\end{center}

部屋の照明を落として、PCの画面光だけで撮った一枚。
\subsection{{\bfseries\sffamily DONE} 不要不急}
\label{sec:orgfe7187d}
人にとって本当に必要な体験というのは、実は「不要不急」の中から生まれるものが多いように思う。

不要不急が制限されているコロナ禍のなかで、人は豊かな生活を求めることができるのだろうか。
\subsection{{\bfseries\sffamily DONE} 貝殻喫茶室}
\label{sec:org998641b}
貝殻弁当をテイクアウトしたときに、店主さんの奥様に「パンもテイクアウトできたら嬉しいですね」と話してみた。

そうしたら、早くも週末に実現させていただけた。パンセットを受け取るときに奥様から「物足りないかもしれません」と言われたけれど、そんなことは全然なかった。

\begin{center}
\includegraphics[width=.9\linewidth]{kaigara-01.jpeg}
\end{center}
\begin{center}
\includegraphics[width=.9\linewidth]{kaigara-02.jpeg}
\end{center}
\subsection{{\bfseries\sffamily DONE} 2021 GW 大瀧詠一}
\label{sec:org901154d}
2021年のGWに入った。今年は4/30も休みにしたので、8日連休になった :-)

大瀧詠一さんを久しぶりに聴きたくなってググっていたら、「ナイアガラ・サウンドのルーツとは？」という記事に出会った。

サイトは文字だけの記事だったので、YouTubeで実際の音源を探してみることにした。

元ネタというか、もう体に染みついているんだろうな。ここは「パクリ」なんて言ってはいけない（笑）
\subsubsection{君は天然色}
\label{sec:org8bf9617}

The Pixies Three - Cold Cold Winter

The Honeycombs - Color Slide
\subsubsection{風立ちぬ}
\label{sec:org3a884b3}
Jimmy Clanton - Venus In Blue Jeans
\subsubsection{夢で逢えたら}
\label{sec:org5f22b6a}
The Shirelles - Foolish Little Girl
\subsubsection{うなずきマーチ}
\label{sec:orgf773528}
The Rivingtons - Papa Oom Mow Mow
\subsection{{\bfseries\sffamily DONE} 2021 GW なろう系小説}
\label{sec:orgb043fb7}
重たい読み物は GW には向いていない。やっぱりライトノベルだな。

話は少し逸れるけれど、ライトウェイトランゲージ（LL）という言葉を今の若い人は知らないらしい。少なくとも、自分より8歳下の職場の人は知らなかった。インタープリタという概念自体は理解してもらえたので、逆にインタープリタ「という呼び方」を知らない人のほうが多いのかもしれないと思うと、ちょっとした違和感がある。
\subsubsection{聖女の魔力は万能です}
\label{sec:org1d8b9c5}
読了。とはいえこの小説はまだ連載中なので、正確には「5月1日時点で読める分まで読了」ということになる。

コミック版もアニメ版もある。というより、なろう系小説はほとんどがそういう形で展開されているものだ。
\subsubsection{転生したらスライムだった件}
\label{sec:org62b67b0}
読了。

こちらは4月時点で読み終えている。アニメのほうは第2シーズンに向けてお休み中。
\subsubsection{無職転生 - 異世界行ったら本気だす -}
\label{sec:orgd3ba1fc}
「なろう系小説」という言葉自体を覚えたのは、実はこの小説がきっかけだった。もっとも、それより前から、なろう系出身のアニメはいくつか見ていたようだ。「転スラ」もそのひとつ。
\subsubsection{これから}
\label{sec:org8851111}
GW はこれからが本番 :laughing:

何を次に読もうかはまだ決めていない。この前届いた「[ダンジョンに出会いを求めるのは間違っているだろうか](\url{https://www.amazon.co.jp/gp/product/B08XM5PBMJ/1})」でも読もうかな。Kindle本はまだ読んでいないものも積み上がっているし :laughing:
ここから下は本文と関係なし。備忘録。

:scream: :boom: :musical\_note: :musical\_note: :hankey:
{[}絵文字チート](\url{https://www.webfx.com/tools/emoji-cheat-sheet/})

TODO<->DONE (C-c C-t / s-<- ->)
{[}Orgチート](\url{https://emacsclub.github.io/html/org\_tutorial.html})
\subsection{{\bfseries\sffamily DONE} ひさしぶりに}
\label{sec:org2a3a64a}
歳を重ねてエンジニアとしての仕事が難しくなり、いわば器用に立ち回ってコンサルタントを生業とするようになってから、そろそろ5年目に入ろうとしている。この仕事もそれなりに面白いのだけど、コミットメントの深さが違うぶん、胸が躍るような刺激までは正直あまりない。

とはいえ、生活のためというほど切実なものでもないし、他に口を出せることもそう多くはない。

ホスピタリティを大事にした対応をすることで、見返りは期待していなくても、育てていく喜びのようなものは感じられている。
\subsection{{\bfseries\sffamily DONE} Yinyoo インイヤーモニター カナル型 CCZ Coffee Bean}
\label{sec:org388a10c}
イヤフォンは、それまで常用していた SHURE SE215（無印）のコネクター部分が壊れてしまったうえに、落として音がシャリシャリになった Apple AirPods Pro もあり、\href{https://www.amazon.co.jp/dp/B09CZHN9VF/ref=sspa\_dk\_detail\_0?pd\_rd\_i=B09CZHN9VF\&pd\_rd\_w=rUZbI\&content-id=amzn1.sym.f293be60-50b7-49bc-95e8-931faf86ed1e\&pf\_rd\_p=f293be60-50b7-49bc-95e8-931faf86ed1e\&pf\_rd\_r=J7R8B10MRFQNY5X13CEZ\&pd\_rd\_wg=mIVsH\&pd\_rd\_r=89b7d26a-db8f-4b66-9e1d-d0d1e9643baf\&s=musical-instruments\&sp\_csd=d2lkZ2V0TmFtZT1zcF9kZXRhaWw\&th=1}{SHURE AONIC 215} もMMXコネクタの接点不良で片耳が聴こえなくなってしまった。

そんな状況の中、中華イヤフォンとしては初めて、昨年（2023年10月）に\href{https://www.amazon.co.jp/1DDダイナミック型HiFiイヤホン-リケーブル可能のデザイン-0-78mm2pinコネクタの採用-バランスな音質表現を持ちハイエンドイヤホン-ミキシング用の中華製イヤホン/dp/B0C39PK494?pd\_rd\_w=BuYlF\&content-id=amzn1.sym.a8cbdd54-57d2-48f2-94b3-7515d6146573\&pf\_rd\_p=a8cbdd54-57d2-48f2-94b3-7515d6146573\&pf\_rd\_r=GH16TVWVRCJ4GN9KA0JJ\&pd\_rd\_wg=rGYXP\&pd\_rd\_r=38b3f7a9-4f01-4a75-86d8-395986f247d3\&pd\_rd\_i=B0C39PK494\&psc=1\&ref\_=pd\_bap\_d\_grid\_rp\_0\_16\_i}{LINKSOULのKiwi Ears Dolce} を購入したところ、有線イヤフォンの奥深い世界にすっかりはまってしまった。

ただ、リケーブルして Dolce を使っていたのだけど、左耳がコネクタから外れて紛失してしまった🥲

その次に買ったのは、同じ KiwiEars の \href{https://www.amazon.co.jp/Linsoul-リケーブル可能のデザイン-0-78mm2pinコネクタの採用-バランスな音質表現を持ちハイエンドイヤホン-ミキシング用の中華製イヤホン/dp/B0C48QLNSQ?pd\_rd\_w=BuYlF\&content-id=amzn1.sym.a8cbdd54-57d2-48f2-94b3-7515d6146573\&pf\_rd\_p=a8cbdd54-57d2-48f2-94b3-7515d6146573\&pf\_rd\_r=GH16TVWVRCJ4GN9KA0JJ\&pd\_rd\_wg=rGYXP\&pd\_rd\_r=38b3f7a9-4f01-4a75-86d8-395986f247d3\&pd\_rd\_i=B0C48QLNSQ\&psc=1\&ref\_=pd\_bap\_d\_grid\_rp\_0\_17\_i}{Linsoul Kiwi Ears Quartet} だったのだけど、サイズが大きくステムも太めで、常用するには向いていなかった。

そこで、ネーミングとコストの両方に惹かれて Amazon でポチったのが \href{https://www.amazon.co.jp/dp/B097XTBYV8?linkCode=sl1\&tag=ymr1107-22\&linkId=01d6a9862f87a25980ed8c8ccafc40b0\&language=ja\_JP\&ref\_=as\_li\_ss\_tl\&th=1}{Yinyoo CCZ Coffee Bean} 👇

寝ながら使うイヤフォンとしても、音質・音圧ともに好みに合っていて、今では一軍として活躍している。

\begin{center}
\includegraphics[width=.9\linewidth]{coffee-bean.png}
\end{center}
\subsection{{\bfseries\sffamily DONE} 葬送のフリーレン}
\label{sec:org795641e}
「葬送のフリーレン」は日本の漫画で、テレビアニメ版が2023年11月から2024年3月にかけて放映された。漫画のほうはまだ連載中。

最初は正直よくわからなかったけど、見進めるうちにどんどん引き込まれる、本当に面白い作品だった。
\subsubsection{アニメ最終話}
\label{sec:orgfa033b5}

アニメのほうはこれで終わった。続きはあるのだろうか🤔

オイサーストで開催された一級魔法使いの試験で出会った、不合格になった受験生との別れのあと、フェルンとシュタルクから「さっぱりした別れ方だね」と言われる。フリーレンが遠い目をしていたその先には、かつての勇者一行、ヒンメルの言葉が思い出されていた。このシーンにこそ、このアニメが伝えたかったことが詰まっている気がする。

また会ったときに、恥ずかしい思いをしたくないからね。

\href{https://watch.amazon.co.jp/detail?gti=amzn1.dv.gti.c74f9d3b-74ad-4853-b591-742901681ec0\&territory=JP\&ref\_=share\_mac\_season\&r=web}{アニメ Eps.28}
\subsection{{\bfseries\sffamily DONE} このサイトのデプロイについて}
\label{sec:org0bbad39}
世の中にはたくさんのブログサービスがあって、そのお手軽さもそれはそれで魅力的だけど、静的サイトジェネレータには静的サイトジェネレータならではの良さがある。動的サイトのようにレンダリング処理が要らないぶん、表示が速いのもポイント。メモ帳程度の文章をつらつら書くだけなので、あまり凝る必要もない。そんな理由から静的サイトジェネレータを選んでいる。

ドメインは自分で取得したものだけど、サイト自体はあちこち渡り歩いてきた。いまは GitHub Pages を使っている。これまでは作ったコンテンツをローカルの\href{https://github.com/ac1965/ac1965.github.io}{GitHub Pages}のリポジトリに \texttt{rsync} してから \texttt{git push} していたのだけど、それすら面倒に感じるくらいのオッさんになってきたので、Post を作って \texttt{git push} したら自動でデプロイされる方向に切り替えたいと思い始めた。

以下、git アカウントの作成や Hugo サイトの構築、git まわりの設定はすでに済んでいる前提で進める。
\subsubsection{GitHub Pagesでデプロイ}
\label{sec:orgc1d3998}
\begin{itemize}
\item \href{https://docs.github.com/ja/pages/getting-started-with-github-pages/about-github-pages}{GitHub Pages について}
\item \href{https://pages.github.com}{GitHub Pages}
\item \href{https://github.com/actions/deploy-pages}{deploy-page}
\end{itemize}

これによると、 \textbf{GitHub Pages} には2つのタイプがあるらしい。

\begin{itemize}
\item User or Organization Site
\texttt{https://<username or organization>.github.io/} のURLで公開されるタイプ
\item Project Site
\texttt{https://<username>.github.io/repository/} のURLで公開されるタイプ
\end{itemize}

リポジトリ名が \texttt{username.github.io} 以外の場合は後者になるということがわかった。たぶん、その理解で合っていると思う。

つまり \url{https://gohugo.io/hosting-and-deployment/hosting-on-github/} のやり方をそのままやると、\url{https://www.ty07.net/hugo-mysite/} のようなURLで公開されてしまう。期待どおり \url{https://www.ty07.net} のURLで公開しようとすると、結局 \href{https://github.com/ac1965/ac1965.github.io}{GitHub Pages} のリポジトリに \texttt{rsync} してから \texttt{git push} する方法しかないのだろうか😓
\subsubsection{すこしたった後に思ったこと。}
\label{sec:orgc3dadf9}

\href{https://github.com/peaceiris/actions-hugo}{actions-hugo} を先に読んでいれば、もう少し理解が早かったかもしれない。要するに \texttt{https://<username>.github.io/} 側に Hugo の環境を入れて、GitHub Pages にデプロイすればいいということがわかった。

\begin{center}
\includegraphics[width=.9\linewidth]{githubpages-setting.png}
\end{center}

\begin{center}
\includegraphics[width=.9\linewidth]{gitubpages-deplogy.png}
\end{center}

ただ、CNAME がうまく処理されなかったので、結局は以前とほぼ同じデプロイ方法に落ち着いた。
\subsection{{\bfseries\sffamily DONE} ThemaをBlowfishにした}
\label{sec:orgdc8c666}
表現したかったことは、たぶん全部実現できた 🤫
\subsection{{\bfseries\sffamily DONE} ひさしぶりに Miles をききたくなった}
\label{sec:org3913891}
\subsubsection{いつまでも憧れ}
\label{sec:org507cc28}
いま58歳で、あと4ヶ月経てば59歳になる。初めて聴いたジャズのプレイヤーが Miles だった。

いつ頃だったかな。長兄が高校に入った頃だから、自分は小学5年生くらいだったと思う。当時はオートリターン機能もついていないレコードプレーヤーで聴いていて、長兄は聴きながらよく寝落ちしてしまうので、演奏が終わったらレコード針を上げるのが自分の役目だった😭

そんなことを思い出しながら \texttt{So What} を聴いている。

日本のCMに出演していたのはいつ頃だったか。VANとTDKのCMがあったな。VANのCMはとにかくかっこよくて、真似したりもした🤣　TDKのほうはうろ覚えなのだけど、亡くなる少し前くらいだったと思う。

Bitches Brew のジャケットアートワークなんかも好きだった。

高校生になる頃からは他のジャズプレイヤーや違うジャンルの音楽を聴くようになっていて、Miles が晩年にファンクやロックを取り入れたアルバムは当時は買わなかった。その代わり、最近はYouTubeやサブスク\textsuperscript{\ref{org3132e76}}で改めて聴き直している。
\subsection{{\bfseries\sffamily DONE} 犬山祭}
\label{sec:orgeccf322}
犬山祭は、愛知県犬山市で毎年4月の第1土曜・日曜に行われる春季祭礼で、かなり大規模なお祭りだ。

犬山には叔父が住んでいるのだけど、58歳になる今日まで、祭にも叔父の家にも一度も行ったことがなかった。それが、ふと思うところがあって、急に行ってみることにした。

叔父は、理由は聞いていないけれど会社員を早期退職して\href{https://gallerykiya.jimdofree.com/kiya工房/}{大工さん}になった人で、自分にとっては父親と同じくらい、生き方や考え方を尊敬している人だ。兄弟だから、どこか似ているところがあるのかもしれない。父は生活のために会社員を続けていたけれど、50代のはじめ（53歳くらいだったか）に、書道家と家業の禰宜を務めるために早期退職している。

犬山祭の試楽（1日目の土曜日）は雨になることが多いらしいのだけど、今年は好天に恵まれて、にぎわいをしっかり感じることができた。祭のほうは用事があったので、車山\textsuperscript{\ref{org3132e76}}の提灯に火が灯った頃に叔父の家に戻り、お茶を飲みながら叔父や従姉妹と少し話をしてからおいとました。
\subsection{{\bfseries\sffamily DONE} 烏は主を選ばない}
\label{sec:org939915c}
阿部智里さんの和風ファンタジー「八咫烏シリーズ」第二弾を原作に、4月6日からNHKでアニメ化された作品。

人間の姿に変身できる一族が、異世界・山内を舞台に繰り広げる物語。魅力的な登場人物たち、周到に張られた謎、迫力あるバトルなど、壮大な世界観が魅力的な作品だ。

原作小説のほうは文藝春秋のウェブサイトで、㊗アニメ化を記念して第一章が全文無料公開されている。
\subsection{{\bfseries\sffamily DONE} 五月晴れの昼下がり}
\label{sec:org20eb316}
GWも終盤。作り置きの焼きそばに追いソースをかけて、お腹が少し満たされたので散歩に出たものの、突き刺すような日差しに負けて家でまったり🥲

洗濯物も気持ちよく乾いていたので、お天道様には感謝。
\subsection{{\bfseries\sffamily DONE} emacs:libgccjit:'libgcc\_s.1.1.dylib' no such file}
\label{sec:org48755d9}
emacsのコンパイルでエラーが出ているのだけど、これは一体どうしたことか。

native-compilation\textsuperscript{\ref{org3132e76}} は libgccjit を使っているのだけど、libgcc\_s.1.1.dylib は gcc の動的ライブラリのはず。

寝ている間に小人さんが直してくれないかな🤔

\begin{verbatim}
otool -L libgccjit.0.dylib
libgccjit.0.dylib:
        /usr/local/opt/libgccjit/lib/gcc/current/libgccjit.0.dylib (compatibility version 0.0.0, current version 26.0.26)
        /usr/lib/libiconv.2.dylib (compatibility version 7.0.0, current version 7.0.0)
        /usr/local/opt/isl/lib/libisl.23.dylib (compatibility version 27.0.0, current version 27.0.0)
        /usr/local/opt/libmpc/lib/libmpc.3.dylib (compatibility  version 7.0.0, current version 7.1.0)
        /usr/local/opt/mpfr/lib/libmpfr.6.dylib (compatibility version 9.0.0, current version 9.1.0)
        /usr/local/opt/gmp/lib/libgmp.10.dylib (compatibility version 16.0.0, current version 16.0.0)
        /usr/lib/libz.1.dylib (compatibility version 1.0.0, current version 1.2.12)
        /usr/local/opt/zstd/lib/libzstd.1.dylib (compatibility version 1.0.0, current version 1.5.6)
        @rpath/libgcc_s.1.1.dylib (compatibility version 1.0.0, current version 1.1.0)
        /usr/lib/libSystem.B.dylib (compatibility version 1.0.0, current version 1345.100.2)
\end{verbatim}
\subsection{{\bfseries\sffamily DONE} ボクが満足するような最強のエディタを作って}
\label{sec:orgb7d0a16}
ChatGPT に聞いてみた。

\begin{quote}
君のことを少しずつ知り始めたばかりだけど、最強のエディタを作るためにいくつか質問させてね。これが君の理想に近づくための第一歩だよ！
\end{quote}
\subsection{{\bfseries\sffamily DONE} 麦わら帽子は}
\label{sec:orgb0633f8}
今年の夏休みも終わった。

吉田拓郎の「夏休み」の歌詞が頭の中で繰り返し流れている。1日の半分くらいは横になって過ごしていた。疲れが溜まっている気分だったので、まあ仕方ない。
\subsection{{\bfseries\sffamily DONE} macOSメモ}
\label{sec:org7b20b2a}
最近はGoogle検索の代わりに \href{https://www.perplexity.ai}{Perplexity AI} を使うことが増えている。調べものの一発目をPerplexityに投げて、参照元をたどりながら裏を取る、という使い方が定着しつつある感じだ。

今回はその流れで、「macOSのシステム設定を移行するコマンドとその手順」をPerplexityに質問してみた。=defaults= コマンド自体は前から名前だけは知っていたんだけど、ドメインの調べ方からエクスポート・インポートまでひととおりまとまった形で聞いたのは初めてだったので、備忘としてそのまま載せておく。

ざっと聞いた感じ、GUIでちまちま設定し直すよりは断然速そうだ。ただ \texttt{defaults export} はドメイン単位でしか吐き出せないので、乗り換えのたびに何のドメインをエクスポートすべきか調べるのが地味に面倒そうではある。とりあえずこの記事をメモ代わりに残しておいて、次に環境を作り直すタイミングで実際に試してみようと思う。
\subsection{{\bfseries\sffamily DONE} EmacsでLLM}
\label{sec:org715d4e7}
昨日、久しぶりに本屋に立ち寄って立ち読みしていたら、MacのOllamaでLLMをローカルで動かして遊ぶという記事を見かけた。普段はあまり技術書コーナーで長居しないんだけど、最近あちこちで「LLMがローカルで動く」という話を耳にしていたので、つい足が止まった。

思わず食指が動いて、さっそくバイナリを落として動かしてみた。手元のマシンでもそれなりにサクサク動いて、これはなかなか遊べそうだったので、年末年始にじっくり勉強することにした。

\href{https://ollama.com/library}{Ollamaで利用可能なモデル}も豊富で、コミュニティが精力的にアップし続けてくれてるみたいだった。これだけ盛り上がっているなら、Emacsで動くパッケージもきっとあるだろうと思って、最近お気に入りの\href{https://www.perplexity.ai}{Perplexity}で調べてみた。\href{https://www.google.com}{Google}でもよかったんだけど、Perplexityのほうがサッと調べてくれるので、最近はこっちに頼りがちだ。

以下、Perplexityの回答をそのまま引用する。

私は\href{https://github.com/conao3/leaf.el}{ leaf }を使っているので、それに合わせて書き換えたものを貼っておく。

\begin{verbatim}
(leaf ellama
  :after llm-ollama
  :ensure t
  :init
  (setopt ellama-language "Japanese")
  (setopt ellama-sessions-directory (concat no-littering-var-directory "ellama-sessions"))
  (setopt ellama-naming-scheme 'ellama-generate-name-by-llm)
  (setopt ellama-provider
          (make-llm-ollama
           ;; this model should be pulled to use it
           ;; value should be the same as you print in terminal during pull
           :chat-model "llama3:8b-instruct-q8_0"
           :embedding-model "nomic-embed-text"
           :default-chat-non-standard-params '(("num_ctx" . 8192))))
  (setopt ellama-summarization-provider
          (make-llm-ollama
           :chat-model "qwen2.5:3b"
           :embedding-model "nomic-embed-text"
           :default-chat-non-standard-params '(("num_ctx" . 32768))))
  ;; llm providers for interactive switching.
  (setopt ellama-providers
          '(("zephyr" . (make-llm-ollama
                             :chat-model "zephyr:7b-beta-q6_K"
                             :embedding-model "zephyr:7b-beta-q6_K"))
            ("mistral" . (make-llm-ollama
                              :chat-model "mistral:7b-instruct-v0.2-q6_K"
                              :embedding-model "mistral:7b-instruct-v0.2-q6_K"))
            ("mixtral" . (make-llm-ollama
                              :chat-model "mixtral:8x7b-instruct-v0.1-q3_K_M-4k"
                              :embedding-model "mixtral:8x7b-instruct-v0.1-q3_K_M-4k"))
            ("codestral" . (make-llm-ollama
                            :chat-model "codestral:22b-v0.1-q4_K_S"
                            :embedding-model "codestral:22b-v0.1-q4_K_S"))
            ("gemma2" . (make-llm-ollama
                         :chat-model "gemma2:27b-instruct-q4_K_S"
                         :embedding-model "gemma2:27b-instruct-q4_K_S"))
            ("command-r" . (make-llm-ollama
                            :chat-model "command-r:35b"
                            :embedding-model "command-r:35b"))
            ("llama3.1" . (make-llm-ollama
                           :chat-model "llama3.1:8b"
                           :embedding-model "llama3.1:8b"))))
  ;; Naming new sessions with llm
  (setopt ellama-naming-provider
          (make-llm-ollama
           :chat-model "llama3:8b-instruct-q8_0"
           :embedding-model "nomic-embed-text"
           :default-chat-non-standard-params '(("stop" . ("\n")))))
  ;; Translation llm provider
  (setopt ellama-translation-provider
          (make-llm-ollama
           :chat-model "qwen2.5:3b"
           :embedding-model "nomic-embed-text"
           :default-chat-non-standard-params
           '(("num_ctx" . 32768)))))
\end{verbatim}
\subsubsection{設定ファイル（Modelfile）の作成}
\label{sec:orgc8f7cae}

\begin{quote}
Ollamaでは、カスタムモデルの設定ファイルとして \uline{Modelfile} が必要です。以下はその例です。

\begin{enumerate}
\item \textbf{Modelfileの作成}:
モデルと同じディレクトリに \texttt{Modelfile} という名前のファイルを作成し、以下の内容を記述します。
\end{enumerate}
\begin{verbatim}
from /ollama/gemma2-9b-chat-gguf
\end{verbatim}
\end{quote}
\subsubsection{Ellamaで遊んだ}
\label{sec:org9dfb271}

\texttt{M-x ellama-chat} で試しに遊んでみたら、思っていたよりずっとちゃんとしていて驚いた🤣

せっかくなので、Emacsの設定ファイルについて聞いてみた。返ってきたコードはあちこち構文が壊れているんだけど、それっぽいコードをそれっぽく組み立ててくる時点で、素直におもしろい。

教科書的な説明を淡々と返してくるだけかと思いきや、ちゃんと動く（つもりの）コードまで組み立ててくるのは、正直ちょっと感動した。エラーが残っている部分はご愛嬌として、年末年始はこのあたりをもう少し触って、実用に耐えるくらいまで設定を詰めていきたい。
\subsection{{\bfseries\sffamily DONE} 備忘録(Emacsのショートカットキー)}
\label{sec:org904f69a}
\subsubsection{機能別キーバインド}
\label{sec:orgd5e7649}
\begin{enumerate}
\item 基本編集操作
\label{sec:org70be987}
\begin{center}
\begin{tabular}{ll}
キーバインド & 機能\\
\hline
C-<backspace> & カーソルの前の文字を削除\\
C-? & ヘルプコマンドを開く\\
C-x u & 元に戻す（undo-fu-only-undo）\\
C-z & やり直す（undo-fu-only-redo）\\
\end{tabular}
\end{center}
\item ウィンドウおよびフレーム管理
\label{sec:org8e58915}
\begin{center}
\begin{tabular}{ll}
キーバインド & 機能\\
\hline
C-. & 他のウィンドウに移動\\
s-o & 他のフレームでファイルを開く\\
s-m & 新しいフレームを作成\\
s-w & 現在のフレームを削除\\
s-. & ウィンドウ分割を切り替える\\
s-j & 他のウィンドウでファイルを開く\\
\end{tabular}
\end{center}
\item ファイルおよびバッファ操作
\label{sec:orgfbdea40}
\begin{center}
\begin{tabular}{ll}
キーバインド & 機能\\
\hline
C-c o & ファイルを開く\\
C-c v & ファイルを読み取り専用モードで開く\\
C-c k & バッファを閉じてウィンドウを閉じる\\
\end{tabular}
\end{center}
\item 検索および置換
\label{sec:org4764bd8}
\begin{center}
\begin{tabular}{ll}
キーバインド & 機能\\
\hline
C-c r & バッファ内の文字列を置換\\
C-c C-r & consult-ripgrepで検索\\
C-s & consult-lineでバッファ内検索\\
M-g g & 行に移動（consult-goto-line）\\
M-g i & バッファ内の関数検索（consult-imenu）\\
M-g b & バッファ切り替え（consult-buffer）\\
\end{tabular}
\end{center}
\item コメントおよび整列
\label{sec:org8e649c8}
\begin{center}
\begin{tabular}{ll}
キーバインド & 機能\\
\hline
C-c ; & 選択範囲をコメントアウト\\
C-c : & 選択範囲のコメントアウトを解除\\
C-c M-a & 正規表現で揃える\\
\end{tabular}
\end{center}
\item スクロールと移動
\label{sec:org88f515b}
\begin{center}
\begin{tabular}{ll}
キーバインド & 機能\\
\hline
C-s-<up> & 上方向にスクロール\\
C-s-<down> & 下方向にスクロール\\
s-<up> & バッファの先頭に移動\\
s-<down> & バッファの末尾に移動\\
s-<wheel-up> & 他のウィンドウをスクロールアップ\\
s-<wheel-down> & 他のウィンドウをスクロールダウン\\
\end{tabular}
\end{center}
\item テキストサイズ調整
\label{sec:org44042df}
\begin{center}
\begin{tabular}{ll}
キーバインド & 機能\\
\hline
C-+ & テキストサイズを増加\\
C-- & テキストサイズを縮小\\
\end{tabular}
\end{center}
\item Orgモード
\label{sec:orgb72509c}
\begin{center}
\begin{tabular}{ll}
キーバインド & 機能\\
\hline
C-c c & Orgエントリをキャプチャ\\
\end{tabular}
\end{center}
\item Emacs制御
\label{sec:org4aed523}
\begin{center}
\begin{tabular}{ll}
キーバインド & 機能\\
\hline
C-q & Emacsを終了\\
M-q & バッファを保存してEmacsを終了\\
s-r & Emacsを再起動\\
\end{tabular}
\end{center}
\item Magit
\label{sec:org8d23764}
\begin{center}
\begin{tabular}{ll}
キーバインド & 機能\\
\hline
C-x g & Magitステータスを開く\\
\end{tabular}
\end{center}
\item Embark
\label{sec:orgdc03265}
\begin{center}
\begin{tabular}{ll}
キーバインド & 機能\\
\hline
s-/ & Embarkアクション\\
s-; & Embark Do What I Mean\\
C-<f2> & Embarkバインディング\\
\end{tabular}
\end{center}
\item その他
\label{sec:orgcfc2308}
\begin{center}
\begin{tabular}{ll}
キーバインド & 機能\\
\hline
M-x & 拡張コマンドを実行\\
\end{tabular}
\end{center}
\end{enumerate}
\subsection{{\bfseries\sffamily DONE} 2026年の生成AI、エンジニア的に気になるトレンド3選}
\label{sec:org210c6b9}
2025年は「AIエージェントの年」なんて言われてた。ChatGPTを筆頭に、各社ツールが職場にわりと浸透してきて、「とりあえず触ってみるか」みたいなノリで実験的に使い始めた企業も多かったと思う。

で、2026年はどうかというと、一言でいえば \textbf{お試し期間終了} って感じ。ROI出せてる？ちゃんと本番で動いてる？って問われる、より実戦的なフェーズに入ってきた。

個人的にもここ最近、生成AI周りのキャッチアップに時間使ってるので、今年気になってるトレンドを3つにまとめてみた。
\subsubsection{その1：AIエージェント、いよいよ本番へ}
\label{sec:orga765221}

これまでの生成AIって、基本「聞いたら答えてくれる」だけだった。プロンプト投げて、テキストが返ってくる。便利だけど受け身な存在。

それが2026年は \textbf{AIエージェント} が本格的に普及してきそう。曖昧な目標を渡すだけで、自律的にワークフローを完遂してくれるやつ。

たとえば「来週の出張手配して」って言うだけで、フライト検索・予算照合・ホテル予約・カレンダー登録まで全部やってくれる、みたいなイメージ。Deloitteは「シリコンベースの労働力」なんて表現してて、AIエージェントを含むハイブリッドな労働力の管理が企業に求められると言ってる。

エンジニア的に押さえておきたいのは、MCP（Model Context Protocol）とかA2A（Agent-to-Agent）あたりの仕様。エージェント同士が連携するための仕組みで、これを理解してないと設計面でちょっと困る場面が増えてきそう。
\subsubsection{その2：マルチモーダル、「扱える」から「わかる」へ}
\label{sec:orga829e6c}

テキスト・画像・音声・動画を横断的に処理するマルチモーダルAIは、2025年にかなり広まった。

2026年はその次のステップで、複数の情報形式を文脈ごとに統合して理解する段階に進んでいる感じ。

具体的なイメージでいうと、会議の録画を放り込んだら、音声だけじゃなくて参加者の表情・スライドの内容・ホワイトボードの手書きメモまで全部解析して、議事録を自動生成してくれる、みたいな話。製造現場での異常音+映像の組み合わせ分析とかも実用化が進んでる。

実装面では引き続きRAGが重要で、ハルシネーション対策の定番テクニックとして定着してきてる。社内データを使ったファインチューニングも、各社で当たり前のように取り組まれてる印象。
\subsubsection{その3：「すごい！」より「で、いくら儲かった？」}
\label{sec:orgdeef02f}

Gartnerによると、2026年末までに世界の企業の80\%以上がGenAI APIやモデルを本格展開するという予測がある。

ただ、もはや「AI導入しました！」だけじゃ評価されない時代で、具体的なコスト削減・収益貢献という数字を出せるかどうかが問われるようになってきた。

エンジニア的にも、技術力だけじゃなくて「このAI、なんでこう判断したの？」を説明できるか、つまりXAI（説明可能なAI）やガバナンス対応が求められる場面が増えてる。

あと、Forresterがディープフェイクの本格的な悪用が2026年に広まると言ってて、検知技術や社内ガイドラインの整備も急ぎたいところ。
\subsubsection{まとめ}
\label{sec:org81bc61d}

\begin{itemize}
\item AIエージェント：設計・統合スキルが差になる
\item マルチモーダル：「扱う」から「統合して理解する」フェーズへ
\item ROI評価：説明責任とガバナンスが求められる時代に
\end{itemize}

「使うかどうか」の議論はもう終わっていて、「いかに使い倒すか」のフェーズに入っている。実装レベルできちんと対応できるかどうかが、今年のエンジニアとしての市場価値に直結してくる気がする。

引き続き、キャッチアップを続けていきたい。
\subsection{{\bfseries\sffamily DONE} .emacs.d を Literate Configuration で管理する}
\label{sec:org61a945c}
\texttt{README.org} を久しぶりに大きく更新したので、設計の話も含めて記録として書き残しておく。
\subsubsection{そもそもなぜこんな構成にしたのか}
\label{sec:org3c0c881}

Emacs の設定は、放っておくと本当にカオスになってしまう。 \texttt{init.el} に \texttt{setq} がだらだらと並んで、どれが何のためにあるのか誰にもわからなくなる（かくいう自分も含めて）。

そこで数年前に、思い切って全部作り直した。コンセプトはシンプルで、

\begin{itemize}
\item 設定をきちんとエンジニアリングの成果物として扱うこと
\item どのモジュールを消しても他が壊れないこと
\item どの環境で起動しても同じ結果になること
\end{itemize}

この3つを最初から保証できる構成にしている。
\subsubsection{Literate Configuration}
\label{sec:org1f592c9}

\texttt{README.org} が唯一の真実の源泉になっている。Emacs Lisp のコードはすべて \texttt{README.org} 内の \texttt{src} ブロックに書いておき、 \texttt{org-babel-tangle} で \texttt{lisp/<layer>/<module>.el} に書き出す仕組みだ。

保存時に自動でタングルされるようにしてあるので、実際の作業フローはほぼ「 \texttt{README.org} を編集して保存するだけ」になっている。

ドキュメントとコードが同じファイルにあるので、「コードはあるけど意図がわからない」という状況が生まれない。これがいちばんの利点だと思っている。
\subsubsection{階層アーキテクチャ}
\label{sec:org0b8a656}

設定は10層に分かれていて、依存関係は必ず下から上への一方向のみ。

\begin{verbatim}
early-init → core → ui → auth → completion → orgx → vcs → dev → utils → personal
\end{verbatim}

\texttt{core} が UI に依存してはいけないし、 \texttt{vcs} が \texttt{dev} に依存してもいけない。この制約を守ることで、モジュールを消したり差し替えたりしても他が壊れないようにしている。

起動時のオーケストレーションは \texttt{modules.el} が担っていて、モジュールごとにエラートラップをかけながらロードしていく。 \texttt{my:modules-verbose} を有効にすると、モジュールごとの起動時間も表示されるので、何が遅いのかも一目でわかる。
\subsubsection{最近の更新（2026-03）}
\label{sec:org4a96cd8}

今月は、わりと大きめの追加が4つあった。
\begin{enumerate}
\item Claude API クライアントを \texttt{utils} 層に追加（Fix AL）
\label{sec:orgbedaa02}

\texttt{utils/utils-claude.el} として、低レベルの HTTP クライアントを作った。公開しているAPIは3つの関数だけ。

\begin{center}
\begin{tabular}{ll}
関数 & 役割\\
\hline
\texttt{utils-claude-get-api-key} & \texttt{auth-source} から API キーを取得\\
\texttt{utils-claude-request-sync} & 同期 POST\\
\texttt{utils-claude-request} & 非同期バリアント（コールバック）\\
\end{tabular}
\end{center}

\texttt{utils} 層に置いたのは、UI にも Org にも依存しない純粋なトランスポート層だから。built-in の \texttt{url} ・ \texttt{json} ・ \texttt{auth-source} しか使っていないので、どの上位層からでも \texttt{require} できる。

\texttt{dev} 層にある \texttt{aidermacs} や \texttt{gptel} とはあくまで別物という整理で、あちらは「開発ツール」、こちらは「HTTP 配管」という位置づけにしている。

エラーは \texttt{utils-claude-error} という専用シンボルで定義していて、HTTP 非200・JSON パース失敗・API キー未設定という3パターンを \texttt{condition-case} でキャッチできるようにしてある。
\item Org Babel から Claude を呼ぶ（Fix AM）
\label{sec:orge38c1df}

\texttt{utils-claude} の上に、個人ワークフロー用のモジュールを \texttt{personal} 層に追加した。

いちばん面白いのは、Org Babel の言語として登録したところで、

\begin{verbatim}
#+begin_src claude :system "You are a code reviewer."
このコードのリファクタリング案を教えてください。
...
#+end_src
\end{verbatim}

これを \texttt{C-c C-c} で評価すると、そのままアシスタントの応答が結果として入る。 \texttt{:system} ・ \texttt{:model} ・ \texttt{:tokens} といったヘッダ引数もサポートしている。

インタラクティブなコマンドも3つ追加した。

\begin{center}
\begin{tabular}{ll}
コマンド & 動作\\
\hline
\texttt{my/claude-refactor-defun} & カーソル位置の \texttt{defun} をリファクタリング依頼\\
\texttt{my/claude-complete-docstring} & \texttt{defun} の docstring を生成\\
\texttt{my/claude-analyze-require-graph} & バッファの \texttt{require} グラフを分析して上向き依存を検出\\
\end{tabular}
\end{center}

\texttt{personal} 層に置いたのは、これが完全に個人的なワークフローの選択だから。共有モジュールのロード順に影響を与えてはいけない、という原則も理由のひとつになっている。
\item ブログ記事キャプチャテンプレート（Fix AN）
\label{sec:org0b3498d}

\texttt{C-c c b} で新しい記事のスケルトンを作れるようにした。

ox-hugo のサブツリーエクスポートは \texttt{* blog} 直下にエントリがないと動かないので、テンプレートのターゲットは \texttt{(file+olp my:f:capture-blog-file "blog")} にしてある。

プロパティドロワーは自動で埋まるようになっている。

\begin{center}
\begin{tabular}{ll}
プロパティ & 値\\
\hline
\texttt{EXPORT\_FILE\_NAME} & UUID（=org-id-new=）\\
\texttt{EXPORT\_DATE} & 現在時刻\\
\texttt{EXPORT\_HUGO\_TAGS} & インタラクティブ入力\\
\texttt{EXPORT\_HUGO\_CATEGORIES} & インタラクティブ入力\\
\texttt{EXPORT\_HUGO\_LASTMOD} & 空（ox-hugo がエクスポート時に補完）\\
\end{tabular}
\end{center}

デフォルトの \texttt{blog.org} は実際には存在しないので、 \texttt{personal/ac1965.el} で \texttt{my:f:capture-blog-file} を \texttt{all-posts.org} にオーバーライドしている。これをやっておかないと、キャプチャ時に \texttt{org-find-olp: Heading not found} というエラーが出てしまう。
\item 保存時に自動で ox-hugo エクスポート（Fix AO）
\label{sec:org9ed0bca}

\texttt{orgx/orgx-auto-hugo.el} を追加して、 \texttt{all-posts.org} を保存したら \texttt{org-hugo-export-wim-to-md :all-subtrees} が自動で走るようにした。

\texttt{orgx-auto-tangle.el} と同じパターンで実装している。

\begin{itemize}
\item \textbf{適格性述語} \texttt{orgx-auto-hugo-{}-{}eligible-p} ： \texttt{file-truename} でシンボリックリンクも解決する
\item \textbf{ガードフラグ} \texttt{defvar-local orgx-auto-hugo-{}-{}in-progress} ：ox-hugo が内部でバッファ保存を再トリガーして無限ループになるのを防ぐ
\item \textbf{ワーカー} \texttt{orgx-auto-hugo-{}-{}maybe} ： \texttt{condition-case} 内で実行する
\item \textbf{フック登録} \texttt{orgx-auto-hugo-{}-{}install-hook} ： \texttt{org-mode-hook} からバッファローカルに \texttt{after-save-hook} を設置する
\end{itemize}

これで「 \texttt{C-c c b} でスケルトン作成 → 書く → 保存」だけで Hugo へのエクスポートまで完結するようになった。この記事自体も、そのフローで書いている。
\end{enumerate}
\subsubsection{まとめ}
\label{sec:orgbafd32d}

今月の更新で、Emacs の中で Claude を使う基盤（ \texttt{utils-claude} ）、それを使った個人ワークフロー（ \texttt{personal-ai} ）、そしてブログ執筆をゼロフリクションにするフロー（キャプチャ＋自動エクスポート）がひととおり揃った。

Literate Configuration ＋ 階層アーキテクチャは、最初に組み上げるコストは高いけれど、長く使い続けるならやってよかったと思っている。「Emacs の設定が育てにくい」と感じている人の参考になれば嬉しい。
\subsection{{\bfseries\sffamily DONE} このブログについて}
\label{sec:org3886a54}
最初はHTMLを手書きしていた。テキストエディタでタグを打って、FTPでアップして、というやつ。今思えばよくやっていたなと思うけど、当時はそれが当たり前だったし、むしろ楽しかった気もする。 \texttt{<html>} とか \texttt{<table>} とか、ちまちまタグを書きながら「これでちゃんと表示されるかな」とブラウザでリロードする、あの感じ。デザインも全部手作りで、我ながら素朴なものだったけど、自分で作った実感があってよかった。

その後、ブログエンジンをいくつか渡り歩いた。PyBlosxom を使っていた時期もあって、Python で書かれた軽量なやつ。テーマやプラグインを自分でいじれるのが気に入っていた。コメント欄のスパム対策や、シンタックスハイライトの設定など、今となってはどうでもいいような細部に妙に時間をかけていた記憶がある。そこからいろいろあって、今は Hugo に落ち着いている。静的サイトジェネレーターなので速いし、Emacs との相性も悪くない。

書いていることは昔から一貫していて、要は「自分が触っているもの」の記録。2009年ごろは Gentoo Linux で Live CD を作っていたり、awesome → xmonad とウィンドウマネージャを渡り歩いていたり。当時の自分は、わりとやんちゃだったなと思う。Guruplug とか EeePC とか、懐かしい単語が記事に並んでいる。あの頃は「小さくて低消費電力の Linux マシン」のようなものに、妙にテンションが上がっていた。今でいう Raspberry Pi みたいな感覚だったのかもしれない。

昔の記事は全部残っているわけではない。手書きHTMLの頃のものは断片的にしか残っておらず、残っているものから起こしている状態。ちゃんとバックアップしておけばよかったとは思うけれど、まあ、それでもいいかなとも思っている。完璧なアーカイブではなく、残っていたものが残っている、くらいの温度感で。むしろ抜けている部分があるからこそ、「あの時期はそういう感じだったんだろうな」と想像する余地があって、それはそれでいい気がしている。

更新頻度は正直低い。気が向いたときや、まとめたいことができたとき、そのくらい。日記代わりにせっせと書くつもりはないし、PV稼ぎにもそれほど興味はない。書きたいから書く、それだけ。下手に「週1更新」なんて決めてしまうと途端につらくなるのが目に見えているので、最初からそういうルールは作っていない。

最近は Emacs の設定まわりや、AI / LLM 関連の話題が多くなってきた。興味が生成AI絡みのことに向くようになってきたので、自然とそちらに引っ張られている感じ。技術トレンドは移り変わりが速いので、自分の中で整理する意味でも、こうして書いておくのは悪くないと思っている。書きながら考えがまとまることも、わりとある。

ちなみに今のワークフローは、Emacs の Org-mode でざっくり書いて、ox-hugo で変換して Hugo でビルドする構成。保存すれば自動でエクスポートされるようになったので、だいぶ楽になった。以前は「書く気はあるのに、投稿するまでの手順が面倒」ということがあったけれど、今はそのあたりのフリクションがほとんどなくなっている。

細々とでも続いているのは、自分のためのメモとして機能しているからだと思う。数年前に書いた設定を見返して「あー、こんなことをやっていたな」となるのが、意外と悪くない。誰かの役に立てばそれはそれで嬉しいけれど、まずは自分のための場所として、これからものんびり続けていくつもり。
\subsection{{\bfseries\sffamily DONE} Claude on Safari}
\label{sec:org96b0a71}
macOS デスクトップアプリ版の Claude で、たまにチャット入力欄がフリーズすることがある。特にシークレットチャットで固まると、履歴が残らないぶん地味に悲しい。

ブラウザ版ならリロードすれば復帰できるので助かるのだが、今度は Safari 特有の別の問題にぶつかった。日本語入力で変換を確定するために押した Enter キーが、そのままメッセージ送信として扱われてしまうのだ。これはこれでストレスが溜まる。
\subsubsection{原因：ブラウザは「確定のEnter」と「送信のEnter」を区別できない}
\label{sec:orga8b3c6c}

調べてみると、原因はブラウザの仕組みにあるらしい。

日本語IMEを使っているとき、変換候補を確定するための「確定Enter」と、メッセージを送るための「送信Enter」は、ブラウザ側からはどちらも単なる \texttt{event.key === 'Enter'} にしか見えない。本来なら \texttt{compositionstart} / \texttt{compositionend} イベントを使って「いま入力中かどうか」を判定すべきところだが、Safari はこのイベントの取りこぼしが Chrome より目立つらしく、確定した直後の Enter がそのまま送信として伝わってしまう。
\subsubsection{対策：Userscripts でこのサイトだけピンポイントに直す}
\label{sec:org0a7689b}

対処法をいくつか検討したが、直したいのは claude.ai だけなので、Safari 拡張機能の Userscripts（開発者：Justin Wasack）でピンポイントに対応することにした。Mac App Store から無料でインストールできるのも決め手のひとつ。

手順は次のとおり。Userscripts をインストールしたら Safari の拡張機能から有効化し、ツールバーのアイコンをクリックして出てくる画面から進める。

貼り付けるコードはこちら。

\begin{verbatim}
// ==UserScript==
// @name         Claude.ai IME Enter Fix (Safari)
// @namespace    local
// @version      1.0
// @match        https://claude.ai/*
// @run-at       document-start
// @grant        none
// ==/UserScript==
(function () {
  'use strict';
  let isComposing = false;
  let lastCompositionEndTime = 0;
  const THRESHOLD_MS = 300;

  document.addEventListener('compositionstart', () => {
    isComposing = true;
  }, true);

  document.addEventListener('compositionend', () => {
    isComposing = false;
    lastCompositionEndTime = Date.now();
  }, true);

  document.addEventListener('keydown', (event) => {
    if (event.key === 'Enter' && !event.shiftKey) {
      const withinGuard = isComposing ||
        (Date.now() - lastCompositionEndTime < THRESHOLD_MS);
      if (withinGuard) {
        event.stopImmediatePropagation();
        event.preventDefault();
      }
    }
  }, true);
})();
\end{verbatim}

保存できたら、実際にちゃんと動くか試してみよう。

これで問題なければ、日本語入力中の誤送信フラストレーションからは解放されるはず :-)
\subsubsection{おまけ：なぜこれで直るのか}
\label{sec:orgba31b98}

気になる人向けに、仕組みをもう少し詳しく書いておく。
\subsection{{\bfseries\sffamily DONE} Emacs 31 は「自分で書いたコード」を減らすチャンスだ}
\label{sec:orgf1e4894}
Emacs 31 の最初のプリテスト版（ \texttt{emacs-31.0.90} ）が公開された。
正式な 31.1 はまだ出ていないが、機能の中身はすでに固まっている。

そこで自分の設定ファイル（ \texttt{\textasciitilde{}/.emacs.d} ）を読み返してみたところ、
ある事実に気づいた。

\textbf{「これ、もう自分で書かなくていいやつだ」というコードが、そこかしこにある。}

長く Emacs を使っていると、設定は自然と分厚くなる。
「本体にない機能」を補うために書いたコードが、少しずつ積み上がっていくからだ。
ところが Emacs 本体も進化する。数年前に自分で書いた機能が、
いつのまにか標準機能として取り込まれていることがある。

そうなったとき、自作コードは「便利な資産」から「余計な負債」に変わる。
本体の実装と微妙にケンカしたり、バージョンアップのたびに壊れたりするからだ。

この記事は、そんな「捨てられる自作コード」を Emacs 31 で棚卸ししてみた記録である。
\subsubsection{その前に：いま Emacs のバージョンはどうなっているのか}
\label{sec:org5e52622}

話を始める前に、事実関係を整理しておく。

\begin{center}
\begin{tabular}{lrl}
呼び名 & バージョン & 状態\\
\hline
安定版 & 30.2 & いま普通に使うならこれ\\
Emacs 31 & 31.0.90 & プリテスト（第1版）。正式リリース前\\
master & 32.0.50 & 開発の最先端。31 の機能をすべて含む\\
\end{tabular}
\end{center}

つまり \textbf{2026年7月時点で Emacs 31.1 はまだ正式リリースされていない} 。
ただし、master ブランチを自分でビルドしている人なら、
31 の新機能はすべて手元で今すぐ試せる。

以下で紹介する機能は、すべて \texttt{etc/NEWS} （Emacs 本体に同梱されている変更履歴）と
実際のソースコードで裏を取ったものだけを挙げている。
「31 で強化されたらしい」という噂レベルの話は入れていない。
\subsubsection{捨てられるもの、その1：tree-sitter まわりの自作コード}
\label{sec:orgce1dca2}

\subsubsection{これまでの苦労}
\label{sec:org5a24a26}

tree-sitter は、プログラムの構造を正確に解析して色分け（シンタックスハイライト）や
インデントを賢くしてくれる仕組みだ。Emacs 29 から使えるようになった。

ただ、これまでは使い始めるまでが面倒だった。

\begin{enumerate}
\item 言語ごとに「文法定義（grammar）」を自分でダウンロード・コンパイルする
\item \texttt{python-mode} を \texttt{python-ts-mode} に読み替える設定を書く
\item 起動時に文法定義が揃っているか確認する
\end{enumerate}

この3つを自動化するために、私は 155 行ほどのコードを書いていた。
\texttt{core-treesit.el} というファイルだ。
\subsubsection{Emacs 31 での変化}
\label{sec:org84e443e}

31 では、この3つが \textbf{ほぼ全部、本体の機能になった} 。

\begin{itemize}
\item \texttt{treesit-enabled-modes} を \texttt{t} にすると、
「文法定義があるモードは自動的に ts-mode 版を使う」ようになる。
モードの読み替え表を自分で書く必要がなくなった。
\item \texttt{treesit-auto-install-grammar} を設定すると、
文法定義がないときに「入れますか？」と聞いて自動導入してくれる。
起動時の一括インストール処理も要らなくなった。
\end{itemize}

結果、155 行あったファイルは \textbf{55 行程度まで縮んだ} 。
残ったのは「どの言語の文法定義をどこから取ってくるか」という一覧表だけである。

\begin{verbatim}
(leaf treesit
  :straight nil
  :if (and (fboundp 'treesit-available-p) (treesit-available-p))
  :custom
  ;; 文法定義がある ts-mode はすべて有効化する
  ((treesit-enabled-modes        . t)
   ;; 文法定義がなければ、確認のうえ自動インストールする
   (treesit-auto-install-grammar . ask)
   (treesit-font-lock-level      . 4)))
\end{verbatim}

たった数行。100 行が消えた。
\subsubsection{ひとつ落とし穴がある}
\label{sec:org62ce095}

\texttt{treesit-enabled-modes} が面倒を見てくれるのは、
\textbf{Emacs に組み込まれている ts-mode だけ} だ。

サードパーティ製のモード（ \texttt{json-mode} や \texttt{toml-mode} など）からの読み替えは対象外なので、
そのままだと期待どおりに動かない。

これは「困った問題」ではなく、むしろ整理のチャンスだと考えている。
\texttt{json-ts-mode} や \texttt{toml-ts-mode} は Emacs 本体に入っているのだから、
サードパーティのパッケージのほうを消してしまえばいい。
パッケージが2つ減る。
\subsubsection{捨てられるもの、その2：「保存時に行末の空白を消す」設定}
\label{sec:org07fe3cd}

これは Emacs 使いなら誰もが一度は書く設定だと思う。

\begin{verbatim}
;; 昔の書き方
(add-hook 'before-save-hook #'delete-trailing-whitespace)
\end{verbatim}

「ファイルを保存するとき、行末の余計な空白を自動的に消す」というものだ。
自分のコードをきれいに保つには良い設定である。
\subsubsection{ただし、これは事故のもとでもある}
\label{sec:orgc9af5e4}

問題は、この設定が \textbf{開いたファイルすべてに効いてしまう} ことだ。

他人が管理しているリポジトリのファイルを、ちょっと確認するために開いて、
うっかり \texttt{C-x C-s} を押したとする。
すると、行末の空白がある全部の行が「変更あり」と判定されてしまう。

たった1文字直すつもりが、200 行の差分が出る。
プルリクエストを送る前に気づけばいいが、気づかないこともある。
これは典型的な事故パターンだ。
\subsubsection{Emacs 31 での解決}
\label{sec:org28aa548}

31 は \texttt{delete-trailing-whitespace-mode} という新しいマイナーモードを用意した。
「全ファイルに効くフック」ではなく、
「そのバッファだけに効くモード」として空白削除を有効にできる。

\begin{verbatim}
(leaf files
  :straight nil
  :custom ((require-final-newline . t)
           (delete-trailing-lines . t))
  ;; プログラム・Org・設定ファイルのときだけ有効にする
  :hook ((prog-mode-hook . delete-trailing-whitespace-mode)
         (org-mode-hook  . delete-trailing-whitespace-mode)
         (conf-mode-hook . delete-trailing-whitespace-mode)))
\end{verbatim}

効く範囲を自分で決められる。これが本来あるべき形だと思う。
\subsubsection{捨てられるもの、その3：「定期的に保存する」ための自作タイマー}
\label{sec:org2904f40}

Emacs には、最近開いたファイル（ \texttt{recentf} ）や、
ファイル内の最後のカーソル位置（ \texttt{saveplace} ）を記録してくれる機能がある。

ところがこれらは、標準では \textbf{Emacs を終了するときにしか保存されない} 。
Emacs がクラッシュしたり、強制終了したりすると、記録が失われる。

そこで多くの人が、こういうコードを書く。

\begin{verbatim}
;; 5分ごとに保存する自作タイマー
(run-with-timer 300 300 #'recentf-save-list)
\end{verbatim}

Emacs 31 は、これを設定項目として用意した。

\begin{verbatim}
(leaf recentf :straight nil
  :custom ((recentf-autosave-interval . 300)))   ; 300秒ごとに自動保存

(leaf saveplace :straight nil
  :custom ((save-place-autosave-interval . 300)))
\end{verbatim}

\textbf{ここで大事な注意点がひとつある。}

これらの設定は「値をセットする」だけでなく、
「タイマーを仕掛ける」という副作用を持っている。
そのため、 \texttt{setq} で書いても \textbf{効かない} 。必ず \texttt{setopt} （leaf なら \texttt{:custom} ）を使うこと。
これは NEWS でもはっきり警告されている。
\subsubsection{捨てられるもの、その4：連打を抑える処理}
\label{sec:org8f2a229}

私の設定には \texttt{orgx} という Org 拡張のレイヤーがあって、
そこで「保存のたびに重い処理が走らないように間引く」コードを3か所書いていた。

たとえば私の設定は README.org という 13,826 行のファイル1本から生成されている。
このファイルを保存するたびに、全体を解析して各設定ファイルを吐き出す（tangle する）。
連続して \texttt{C-x C-s} を押すと、そのたびに 13,826 行の処理が走る。重い。

そこで「最後の保存から2秒経つまで待つ」という制御を自分で書いていた。
プログラミングの世界では「デバウンス」「スロットル」と呼ばれる、よくあるパターンだ。

Emacs 31 は \texttt{timeout.el} というライブラリを標準で同梱するようになった。

\begin{verbatim}
(require 'timeout nil t)

;; 「2秒に1回まで」に制限する
(timeout-throttle #'orgx-auto-tangle--tangle 2.0)
\end{verbatim}

1行である。

ただし、これで自作コードを全部消していいわけではない。
\texttt{timeout-throttle} は「呼び出し頻度を制限する」だけで、
「すでに走っている処理に二重で入ってくるのを防ぐ」わけではない。
再入防止のフラグは \textbf{残しておくべき} だ。

こういう「本体機能で置き換えられる範囲と、そうでない範囲」を見極めるのが、
リファクタリングでいちばん頭を使うところである。
\subsubsection{ここからは「増やす」話：Emacs 31 の新機能で嬉しかったもの}
\label{sec:org8d36dad}

捨てる話ばかりしてきたが、31 には積極的に取り入れたい新機能も多い。
特に印象的だったものを挙げる。
\subsubsection{Git 操作が本体だけでかなり戦えるようになった}
\label{sec:orgc577f8b}

Emacs で Git を使うなら Magit、というのがここ10年の常識だった。
Emacs 標準の VC は、正直そこまで便利ではなかったからだ。

ところが 31 の NEWS を読むと、VC 関連の変更項目が \textbf{40 を超えている} 。

\begin{itemize}
\item 複数 worktree の操作
\item cherry-pick / revert
\item incoming / outgoing の差分表示
\item コミットの非同期実行（大きなリポジトリで待たされない）
\item 外部で git 操作したときのバッファ自動更新
\end{itemize}

とりあえず、事故防止に効くものだけでも入れておく価値がある。

\begin{verbatim}
(leaf vc
  :straight nil
  :custom
  (;; 公開済みの履歴を書き換えようとしたら止める（force push 事故の防止）
   (vc-allow-rewriting-published-history . nil)
   ;; コミットを非同期にする
   (vc-async-checkin                     . t))
  :config
  (vc-auto-revert-mode 1))
\end{verbatim}

「じゃあ Magit を消せるか」というと、そこは急がないほうがいい。
私は \textbf{まず VC 側を有効にして、1〜2か月ふつうに使ってみてから決める} つもりだ。
道具を捨てる判断は、実際に使ってみた実感でしか下せない。
\subsubsection{ターミナルの Emacs でもポップアップが出る}
\label{sec:org420e835}

これは地味だが大きい。

31 では、ターミナル（TTY）上の Emacs でも「子フレーム」が使えるようになった。
これまで GUI 版でしか出せなかった補完のポップアップ（ \texttt{corfu} など）が、
\texttt{emacs -nw} でも動く可能性がある。

一方で、同じ流れで \textbf{ターミナルのマウス操作モードが標準で有効になる} 変更も入った。
これは iTerm2 や Ghostty で「ドラッグして選択してコピー」ができなくなる、
という体感の変化を生む。

ターミナルで Emacs を使う人は、方針を決めておいたほうがいい。

\begin{verbatim}
;; 端末側のコピー&ペーストを優先したい場合
(unless (display-graphic-p) (xterm-mouse-mode -1))
\end{verbatim}
\subsubsection{macOS 向けの新機能がふたつ}
\label{sec:orgfb5cf3f}

私は Mac で Emacs を使っているので、これは嬉しかった。

\begin{itemize}
\item \texttt{system-taskbar-mode} — Dock アイコンにバッジや進捗バーを出せる
\item \texttt{system-sleep} — 長時間処理の最中にスリープするのを防げる
\end{itemize}

たとえば「Emacs 本体を自分でビルドしている最中に Mac が寝てしまう」という
地味な困りごとが、これで解決する。

\begin{verbatim}
(when (and (eq system-type 'darwin) (fboundp 'system-taskbar-mode))
  (system-taskbar-mode 1))
\end{verbatim}

さらに \texttt{native-comp-async-on-battery-power} という設定が増えた。
Emacs は裏でコードを機械語にコンパイルして高速化する仕組み（native-comp）を持っているが、
これがノートパソコンのバッテリーをけっこう食う。
バッテリー駆動中はコンパイルを止める、という指定ができるようになった。

\begin{verbatim}
(setopt native-comp-async-on-battery-power nil)
\end{verbatim}
\subsubsection{コードの折りたたみが実用レベルになった}
\label{sec:org52a6e6c}

Emacs には昔から \texttt{hideshow} という「関数を折りたたむ」機能があるのだが、
正直あまり使い勝手が良くなく、多くの人はサードパーティ製で代替していた。

31 でこれが一気に強化された。
折りたたみの循環（親→入れ子→全展開）、余白への目印表示、
隠れている行数の表示など、必要なものが一通り揃った。

\begin{verbatim}
(leaf hideshow
  :straight nil
  :hook ((prog-mode-hook . hs-minor-mode))
  :custom
  ((hs-show-indicators      . t)     ; 折りたたみ位置に目印を出す
   (hs-display-lines-hidden . t))    ; 何行隠れているか表示する
  :bind ((:hs-minor-mode-map
          ("<backtab>" . hs-cycle))))
\end{verbatim}
\subsubsection{移行するときに気をつけること}
\label{sec:orgd536905}

新しいバージョンに上げるとき、いちばん怖いのは「気づかないうちに壊れている」ことだ。
31 の NEWS には「互換性のない変更」という節があり、
そこから自分の設定に影響しそうなものを拾っておいた。

\begin{center}
\begin{tabular}{ll}
変更点 & どう困るか\\
\hline
\texttt{if-let} / \texttt{when-let} が非推奨に & 末尾に \texttt{*} が付いた形（=if-let*=）に書き換えが必要\\
ts-mode のインデント設定の変数名が変わった & 独自に設定していると効かなくなる\\
Pcase の入れ子バッククォートが使えない & \texttt{pcase} を凝った使い方をしていると壊れる\\
文字列の書き換え制限が厳しくなった & リテラル文字列を破壊的に変更するコードがエラーに\\
Org 9.8 が同梱される & 自分で Org のバージョンを固定していると二重管理になる\\
\texttt{site-start.el} の読み込み順が変わった & Homebrew などが置いたファイルの挙動が変わりうる\\
\end{tabular}
\end{center}

これらは、いきなり本番の設定に適用してハマるのではなく、
\textbf{ブランチを切って、ひとつずつ試す} のが確実だ。

私の場合はこういう手順で進めている。

\begin{enumerate}
\item \texttt{git switch -c feat/emacs31} でブランチを切る
\item 「自作コードを消す」系（tree-sitter、空白削除、タイマー）を先に適用する
\item 消した関数がどこかから呼ばれていないか grep で確認する
\item リスクの低い小さな設定をまとめて適用する
\item 大きめの新機能は \textbf{1項目ずつ} 入れて、そのつど起動時間とメモリを測る
\item 実測してから、不要になったパッケージを消す判断をする
\end{enumerate}

特に 5 が大事だと思っている。
「新機能を全部入れたら遅くなった」となったとき、
まとめて適用していると、どれが原因かわからなくなる。
\subsubsection{結局、何が得られるのか}
\label{sec:org8e319ae}

数字で言えば、こうなる。

\begin{center}
\begin{tabular}{lll}
項目 & 現状 & 移行後\\
\hline
tree-sitter 設定 & 約155行 & 約55行\\
自作タイマー & 手書き & 設定項目1つ\\
自作デバウンス & 3か所 & ライブラリ呼び出し\\
削除できるパッケージ & — & 最大5個\\
\end{tabular}
\end{center}

ただ、正直なところ \textbf{行数が減ること自体はどうでもいい} 。

本当に意味があるのは、
「自分が保守しなければならないコード」が
「Emacs 本体が動作を保証してくれるコード」に置き換わったことだ。

自作コードは、Emacs がバージョンアップするたびに壊れる可能性がある。
本体の機能は、そう簡単には壊れない（壊れるときは NEWS に書いてある）。

設定を育てるというのは、コードを足していくことだと思われがちだが、
実際には \textbf{「本体が追いついたら、自分のコードを引っ込める」} という
引き算の作業も同じくらい重要なのだと、今回あらためて感じた。
\subsubsection{おまけ：自分で確かめる方法}
\label{sec:org74149e8}

この記事の内容は、すべて Emacs 本体のソースと NEWS で確認できる。
噂や又聞きではなく、一次情報にあたる方法を書いておく。

\begin{verbatim}
# Emacs の中から NEWS を読む
C-h n            # view-emacs-news
C-u C-h n        # バージョンを指定して絞り込む

# ソースから直接確認する
git -C ~/src/emacs show emacs-31:etc/NEWS | less
git -C ~/src/emacs grep -n "defcustom treesit-enabled-modes" emacs-31 -- lisp/treesit.el
\end{verbatim}

新機能の正確な引数や挙動は、 \texttt{C-h f} で関数のドキュメントを引くのがいちばん早い。
ブログ記事（この記事も含めて）を鵜呑みにせず、手元で確かめてほしい。
\subsection{{\bfseries\sffamily DONE} GNU Machのソースコードを、素人が興味本位で読んでみた話}
\label{sec:org1c6067b}
きっかけは本当に些細なことで、「GNU Hurdって結局動いてるの?」という
Twitter(もうXって呼ぶべきなんだろうか)のタイムラインの雑談を見かけた
ことだった。Hurdというと「1990年代からずっと開発中と言われ続けているOS」
という、ある種のミームみたいな扱いをされがちだけど、その基盤になっている
マイクロカーネル \texttt{gnumach} の中身を実際に読んだことはなかったな、と
思い立って、素朴に \texttt{git clone} して読み始めた。

半分は職業病みたいなものかもしれない。「動いていると言われているもの」
の中身を実際に確認する、という行為自体にわりと快感を覚えるタイプの人間
だ。とはいえ今回は仕事とは何の関係もない、純粋な趣味の話。

以下、読んだ範囲で「へぇ」と思ったところを、初見の人にも伝わるように
整理してみる。実は25年ほど前にも一度Hurd周りのコードを読んでいた
時期があって、後半はその頃の記憶を辿る話にもなっている。
\subsubsection{そもそも gnumach って何なのか}
\label{sec:orgf7f1fe3}

\href{https://github.com/flavioc/gnumach}{GNU Mach} は、GNU Hurd の土台になっている
マイクロカーネルだ。ルーツを辿ると Carnegie Mellon 大学の Mach 3.0 に
行き着く、由緒正しい(?)1980年代の設計。「マイクロカーネル」という
言葉の教科書的な説明——「カーネルはメカニズムだけ提供し、ポリシーは
ユーザ空間に追い出す」——を、これほど愚直に実装しているコードベースも
珍しいんじゃないかと思う。

ディレクトリ構成を見るだけでも性格がわかる。

\begin{center}
\begin{tabular}{lll}
ディレクトリ & 中身 & 規模\\
\hline
\texttt{kern/} & スケジューラ・タスク/スレッド管理 & 約31,500行\\
\texttt{ipc/} & Mach IPC(ポート・メッセージング) & 約16,600行\\
\texttt{vm/} & 仮想メモリ管理 & 約21,200行\\
\texttt{device/} & デバイス抽象層 & 約9,700行\\
\texttt{linux/} & Linuxから移植したドライバ群 & 約345,500行\\
\end{tabular}
\end{center}

見ての通り \texttt{linux/} が圧倒的に大きいんだけど、これは中身がMachというより
「LinuxのドライバをMachの作法でラップしたもの」がほとんどだからだ。
IDE/SCSI/AHCIのディスクドライバやネットワークカードのドライバを、いちいち
Machのために書き直すのは非現実的なので、既存資産を持ってきている。この
あたりの「潔さ」は好きだ。

肝心の「マイクロカーネルらしさ」が詰まっているのは \texttt{kern/} ・ \texttt{ipc/} ・ \texttt{vm/}
の3つで、合計しても7万行に満たない程度。今回はほぼこの3つだけを読んだ。
\subsubsection{一番おもしろかった:ポートは「番号」でしかない、という設計}
\label{sec:org1e78284}

Mach IPCの根幹は「すべてはポートへのメッセージ送信である」という一文に
尽きる。ファイルもプロセスもメモリオブジェクトも、カーネルの目からは
「ポート」という統一的な窓口を通してやり取りされる。

で、実際にコードを読んでいちばん「なるほど」と思ったのが、ユーザ空間から
見えている「ポート」の正体だ。 \texttt{ipc/ipc\_kmsg.c} の \texttt{ipc\_kmsg\_copyin\_header}
という関数(送信メッセージのヘッダを処理する箇所)を読むと、こう書いて
ある。

\begin{verbatim}
is_read_lock(space);
entry = ipc_entry_lookup(space, dest_name);      // 名前→エントリ
...
dest_port = (ipc_port_t) entry->ie_object;
ip_lock(dest_port);
is_read_unlock(space);
...
dest_port->ip_srights++;                          // 参照カウント増加
\end{verbatim}

つまりユーザ空間が持っている「ポート」というのは、実体としてはただの
整数(\texttt{mach\_port\_name\_t})だ。カーネルオブジェクトへのポインタなんて
どこにも渡っていない。その整数を、必ず「呼び出したタスク自身のIPC
空間」というテーブルに引かせて、初めて実オブジェクトのポインタが手に
入る。

これがどういう意味を持つかというと、*他人のポート番号を知っていても、
それだけでは何もできない* ということだ。番号が指すテーブルは常に
「自分の空間」なので、他タスクの空間を覗く経路がAPIレベルで存在しない。
UNIXのファイルディスクリプタと似ているようで、「送信権」「受信権」を
はっきり区別した、もう一段厳格なケーパビリティモデルになっている。

しかも権利を人に渡すときのコストが、渡し方によってきっちり作り分けられて
いるのも面白いところだ。 \texttt{ipc/ipc\_right.c} の \texttt{ipc\_right\_copyin} を見ると:

\begin{itemize}
\item \texttt{MAKE\_SEND} (受信権から新しい送信権を複製する)は、単に参照カウントを
増やすだけ
\item \texttt{MOVE\_SEND} (自分の送信権を他人に渡して、自分は手放す)は、内部の
「ユーザ参照カウント(urefs)」をちゃんとデクリメントし、0になったら
エントリごと消す
\end{itemize}

「コピーする」と「明け渡す」で明確にコード上の扱いが違う、というのは、
1980年代の設計とは思えないほど几帳面だなと感心した。
\subsubsection{ページフォールトの中身は、ほぼ「シャドウチェーンを遡る」だけ}
\label{sec:org3fec270}

もう一つ面白かったのが仮想メモリ側、 \texttt{vm/vm\_fault.c} のページフォールト
処理だ。

Machのメモリオブジェクトは「シャドウチェーン」という親子関係を持って
いて、あるページが自分のオブジェクトに無ければ親を見に行く、という
仕組みになっている。コードで言うとこの部分だ。

\begin{verbatim}
if (m->absent) {
    offset += object->shadow_offset;
    next_object = object->shadow;
    if (next_object == VM_OBJECT_NULL) {
        /* シャドウチェーンの末端 = ゼロフィルページを生成 */
        real_m = vm_page_grab(VM_PAGE_HIGHMEM);
        ...
        vm_page_zero_fill(m);
    }
    /* else: 親オブジェクトを見に行く */
}
\end{verbatim}

これ、要するに \texttt{fork()} したあとの匿名メモリ(ヒープやBSS)がどう実現
されているか、そのものだ。「無ければ親を見る、それでも無ければゼロで
埋める」——たったこれだけのルールの積み重ねで、プロセスのメモリ空間が
成立している。

Copy-on-Writeの実装も同じ発想の延長で、書き込みが起きた瞬間に「元の
ページの中身を、コピー先のオブジェクトに逃がしてから」元を書き換え
可能にする、という順序を守っているだけだ。派手な仕掛けは何もなくて、
地味なルールの積み重ねだけでOSの根幹機能が成立しているのを見ると、
なんというか、静かに感動する。
\subsubsection{カーネルがポリシーを持たない、を本当に徹底している話}
\label{sec:org06dd3ef}

個人的に一番「これは筋がいい設計だ」と思ったのが、VMのコピー戦略の
決め方だ。

\texttt{vm/vm\_object.c} の \texttt{vm\_object\_copy\_strategically} という関数は、メモリを
コピーするときに「即座に全部コピーする」「遅延COWにする」「ページャに
丸投げする」の3通りをswitchで切り替えているんだけど、この戦略を
\textbf{カーネル自身は決めていない} んだよね。

\begin{verbatim}
while (!src_object->pager_ready) {
    vm_object_wait(src_object, VM_OBJECT_EVENT_PAGER_READY, interruptible);
    ...
}
\end{verbatim}

外部のページャ(ユーザ空間で動くサーバプロセス)が「私が管理するこの
メモリはこの戦略でコピーしてください」と申告してくるのを、カーネルは
律儀に待つだけなのだ。申告が届くまではコピー処理そのものをブロック
する。

Hurd側の対応も追ってみたら、ファイルシステム翻訳者(ext2fsなど)が
使う \texttt{libpager} というライブラリの \texttt{pager\_create()} 関数が、まさに
この \texttt{copy\_strategy} を引数として受け取っていた。実行ファイルや
共有ライブラリのマップに使われる「遅延COW」は、カーネルが勝手に判断
しているのではなく、ファイルシステム側のサーバが「これは遅延COWで
いいですよ」と明示的に宣言しているから成立している、ということだ。

「メカニズムはカーネル、ポリシーはユーザ空間」という設計哲学、教科書
では何度も見た文言だけど、実際にコードでここまで愚直に貫かれていると、
ちょっと感動してしまう。
\subsubsection{余談:25年前に読んでいたのはHurd側だった}
\label{sec:orgcb69d0a}

ここで少し個人的な話をしたい。今回gnumach(カーネル側)を読んで
いて気づいたんだけど、実は25年ぐらい前、一度Hurd/gnumachのソースを
読もうとしたことがある。当時惹かれていたのは
カーネル側ではなく、Hurd側の \textbf{translator} という仕組みだった。

translatorというのは平たく言うと「ファイルシステムのノード1つ1つに、
好きなユーザ空間サーバを貼り付けられる」機構だ。`settrans`コマンド
一発で、あるディレクトリの正体を任意のプログラムにすり替えられる。
2005年前後にFUSEが登場したときに「これは便利だ」と持て囃されたけど、
Hurdはその発想をずっと早く、しかもマウントポイント単位の局所的な話
ではなく、ファイルシステム階層そのものの設計原理としてやっていた。

当時特に読んでいたのは \texttt{pfinet} だ。ネットワークをファイルシステムの
translatorとして実装するという、今思えばかなり大胆な発想だった。今回
gnumachを読み返したついでに、記憶を頼りに \texttt{pfinet} のソースも覗いて
みたところ、READMEにこう書いてあった。

\begin{quote}
"The Hurd's pfinet server is based on networking code taken from the
Linux kernel sources, initially version 2.2.12 of Linux."
\end{quote}

つまり中身は文字通り、 \textbf{Linux 2.2.12のTCP/IPスタックをほぼそのまま移植したもの} だ。
\texttt{pfinet/linux-src/net/\{core,ipv4,ethernet\}} という
ディレクトリがまるごと存在していて、外側の \texttt{socket-ops.c} ・ \texttt{io-ops.c}
がHurdのRPC作法(`io\_read`/`io\_write`、`socket\_create`など)を実装し、
内側でそのままLinuxのソケットコードを叩く、という二層構造。前段で
書いた「translatorは`memory\_object`の発想の応用である」という話で
言うと、pfinetは「ネットワークスタックという巨大な機能ブロックを丸ごと
ユーザ空間サーバに閉じ込めた」かなり極端な例だったんだと思う。
カーネルはネットワークのことを一切知らない。

当時は結局、自分でtranslatorを書くところまでは至らず、読むだけで
終わった。`trivfs`まわりの作法(`trivfs\_S\_dir\_lookup`のあたり)を
理解しようとしたところで力尽きた記憶がある。それでも、25年経って
また同じコードの子孫を読み返すことになるとは思っていなかったので、
不思議な巡り合わせだった。

ちなみに、25年前の設計が令和のいまも現役だという証拠に、2026年Q2の
Hurd開発ニュースにもちょうどpfinetの話が出ていた。 \texttt{mmap} の戻り値
をチェックせずに使っていた古いバグが2つ見つかり、メモリが逼迫すると
pfinetがクラッシュしうる、という修正が最近入ったばかりだったそうだ。
1999年頃のコードが、2026年になっても現役でバグ報告と修正を受けている
というのは、なかなか感慨深いものがある。
\subsubsection{aarch64移植、見た目と実態のギャップが面白かった}
\label{sec:orgaa5cda7}

ここからは余談。gnumachには \texttt{aarch64/} というディレクトリがあって、
「お、ARM移植も進んでるのか」と覗いてみたら、中身は型定義とABIの
ヘッダファイルだけで、 \texttt{.c} ファイルが1つも無い。ブートコードも例外処理
コードも見当たらない。

「あれ、これって看板倒れなのでは……」と思ってコミットログを辿ったら、
直近のコミットメッセージにこんな一文があった。

\begin{quote}
"Bugaev's wip-aarch64 worked around this by editing kern/lock.h itself\ldots{}"
\end{quote}

つまり、実際に動くARM移植は別ブランチ(\texttt{github.com/bugaevc/gnumach} の
\texttt{wip-aarch64})に存在していて、mainリポジトリには型定義だけが先行して
着地している状態だった。少し検索してみたところ、開発者本人が

\begin{quote}
"my branch of gnumach compiles, boots in QEMU, and passes the testsuite!"
\end{quote}

と、2024年の時点でQEMU上でブートしてテストスイートまで通していると
報告していて、fork()やシグナル配送(intr RPC → シグナルトランポリン →
sigreturn)まで一通り動く状態になっているとのこと。2026年に入ってからは、
別の開発者がこの動くブランチを「共有ファイルを一切壊さずにmainへ
段階的に取り込む」という、地味だけど誠実な作業を進めているようだ。

コミットログという「表の記録」だけを見ていると進捗を過小評価してしまう
典型例だったな、と反省した。オープンソースのアーキ移植って、動く
実装があってもupstream化には別の労力がかかるんだな、というのを実感
した瞬間だった。
\subsubsection{直近のバグ修正から見える、地味だけど大事な話}
\label{sec:orgd9db471}

最後にもう一つ、コミット履歴を追っていて印象に残ったバグ修正を紹介
する。

ひとつは、SMP環境でスケジューラがNULLポインタを踏んでカーネルパニック
していたバグ。コミットメッセージには「Hurdを自己ホストビルドしている
最中にSMP環境でクラッシュする」とあって、地味に実害の大きいバグだった
んだろうなと想像がつく。

もうひとつは、もう少し本質的な話。 \texttt{task\_priority()} という
RPC——タスク自身が自分の優先度を変更するための呼び出しなんだけど、特権
ポートを一切要求しない、つまり \textbf{どのプロセスでも呼べる} ものだ。
これに、修正前は上限チェックが一切なかった。

\begin{verbatim}
task_lock(task);
if (task->max_priority > priority) {      // ← 追加されたチェック
    task_unlock(task);
    return KERN_NO_ACCESS;
}
task->priority = priority;
\end{verbatim}

Machの優先度は数値が小さいほど「偉い」扱いなので、これはつまり「誰でも
自分を最優先に格上げできてしまい、他のタスクを飢餓状態に追い込める」
状態だったということだ。修正では \texttt{max\_priority} という天井を導入し、
それを超える要求は拒否するようにしていた。

CVE番号が付くような華々しい脆弱性ではなく、開発者自身が地道なコード
レビューで見つけて塞いだ類の話だと思うけど、「マイクロカーネルに
おいて、タスクポートは無制限の自己管理権限ではない」という信頼境界を
守る、というのはまさにこのOSの根幹に関わる話だよなと感じた。
\subsubsection{まとめ}
\label{sec:org001eee0}

読んでみて改めて思ったのは、gnumachというコードベースは全体としては
「枯れた設計」なんだけど、枯れているからこそ設計思想がむき出しで見える、
という面白さがあるということだった。ポートは番号でしかない、コピー
戦略はユーザ空間が決める、カーネルは律儀に待つ——1980年代に確立された
アイデアが、令和のいま2026年のコミットログの中でも変わらず一貫して
守られているのを見るのは、なかなか得難い体験だった。

気が向いたら、スケジューラ(\texttt{sched\_prim.c})のロードバランシング周りも
続けて読んでみようと思う。
\subsection{{\bfseries\sffamily DONE} evalの中身を覗いてみた ─ Emacsが40年抱えてきた矛盾の話}
\label{sec:org85d579b}
普段は \texttt{README.org} をいじって設定を壊しては直す、という生活をしているが、たまには自分が毎日呼び出している \texttt{eval} が実際どう動いているのか見てみたくなった。というわけでGNU Emacs本体のCソース、 \texttt{src/eval.c} (4,700行弱)を読んでみた記録。読み終えて分かったのは、この評価器が「言語処理系としての正しさ」より「1985年から動いているコードを1行も壊さないこと」を一貫して優先し続けてきた、継ぎ足しの塊だということだった。
\subsubsection{eval\_subという名の交通整理係}
\label{sec:orgfb634b2}

Emacs Lispで何かを評価するとき、内部では \texttt{eval\_sub} という関数がすべてのS式を捌いている。中身は驚くほど単純な分岐だ。

\begin{itemize}
\item コンスセルでなければ、そのまま返す(数値や文字列の「自己評価」)
\item シンボルなら、変数として値を引く
\item コンスセルなら、関数呼び出しか特殊形式かをさらに判定する
\end{itemize}

コンスセルの場合の分岐がこの関数の本体で、先頭要素(関数名にあたる部分)の型によって以下のように処理が分かれる。

\begin{itemize}
\item Cで書かれたプリミティブ(SUBR)ならC関数を直接呼ぶ
\item レキシカルクロージャやネイティブコンパイル関数なら \texttt{apply\_lambda} に委譲
\item \texttt{(macro . 関数)} という形のコンスセルならマクロとして展開してから再評価
\item \texttt{(lambda ...)} という生のコンスセルなら、これも \texttt{apply\_lambda} に委譲
\end{itemize}

面白いのは、Emacs Lispの「関数」が実際には少なくとも4通りの異なる内部表現(Cのポインタ、コンパイル済みクロージャ、生のconsラムダ、ネイティブコード)を持っているのに、呼び出し側からはその違いが完全に隠蔽されている点だ。標準関数を後からユーザーが純Lispで再定義しても、呼び出し規約はまったく同じまま動く。
\subsubsection{特殊形式は「引数を評価しないC関数」でしかない}
\label{sec:org0601ad0}

\texttt{if} や \texttt{let} や \texttt{while} といった特殊形式は、 \texttt{UNEVALLED} という印がついたSUBRとして定義されている。

\begin{verbatim}
DEFUN ("if", Fif, Sif, 2, UNEVALLED, 0, ...)
  (Lisp_Object args)
{
  Lisp_Object cond = eval_sub (XCAR (args));
  if (!NILP (cond))
    return eval_sub (Fcar (XCDR (args)));
  return Fprogn (Fcdr (XCDR (args)));
}
\end{verbatim}

つまり \texttt{(if cond then else)} と書いたとき、Emacsは \texttt{cond} の部分だけを評価してから、Cのif文で \texttt{then} か \texttt{else} のどちらを評価するかを自分で決めている。マクロ展開という別フェーズがあるわけではなく、評価器がその場その場で「これは何を評価すべきか」を判断しているだけ、というのが素朴で好きなところだ。ちなみに \texttt{if} の内部で \texttt{else} 節の評価を丸ごと \texttt{Fprogn} に投げているのを見ると、「 \texttt{if} は \texttt{progn} に依存して実装されている」という実装同士の依存関係も透けて見える。
\subsubsection{nilが自分自身を返す理由}
\label{sec:org917ba67}

\texttt{nil} を評価すると \texttt{nil} が返るのは当たり前に思えるが、実はシンボルなので本来は「変数として値を引く」経路を通るはずだ。種を明かすと、 \texttt{nil} ・ \texttt{t} ・キーワードシンボルは初期化時に「自分自身を値として持つ」よう仕込まれた上で、書き換え禁止フラグ(\texttt{SYMBOL\_NOWRITE})が立てられている。 \texttt{data.c} の \texttt{set\_internal} を見ると、このフラグが立ったシンボルへの再代入は問答無用でエラーになる(キーワードが自分自身に再設定される場合だけ例外)。

見た目の「自己評価」と、C言語レベルでの実装はまったく別物というのが地味に面白かった。アトムの自己評価は「コンスでなければスルーする」だけの話で、 \texttt{nil} / \texttt{t} /キーワードの自己評価は「シンボルの値セルが定数として固定されている」という、まったく別の仕組みで実現されている。
\subsubsection{レキシカルスコープを動的束縛の上に建て増しした話}
\label{sec:orgf9cbb96}

一番驚いたのはここだ。Emacs Lispは今でこそレキシカルスコープ(クロージャがちゃんと変数を閉じ込める方式)を使えるが、これは後付けの機能で、しかも \textbf{内部実装はダイナミックスコープの仕組みをそのまま流用} している。

具体的には、レキシカル環境そのものが \texttt{Vinternal\_interpreter\_environment} という、特殊変数(つまりダイナミックに束縛される変数)の中にalistとして格納されている。 \texttt{let} でレキシカル変数を作るというのは、実は「環境alistというダイナミック変数に新しいコンスセルをconsして、それをまたダイナミックに束縛し直す」という操作に過ぎない。土台はダイナミックスコープのまま、その上にレキシカルスコープごっこを乗せているような設計だ。

そして、シンボルごとに「動的」か「レキシカル」かを切り替えるのが \texttt{declared\_special} フラグだ。これは \texttt{(defvar sym val)} のように初期値つきで宣言すると立つ、シンボルに一生ついて回るグローバルな印になる。

さらに面白いのが \texttt{(defvar foo)} (初期値なし)の挙動だ。これは、そのシンボルを「このスコープでは動的扱いにする」という \textbf{一時的な} 目印として、環境alistに \textbf{コンスセルではなく裸のシンボル} を突っ込む。

\begin{verbatim}
Vinternal_interpreter_environment
  = Fcons (sym, Vinternal_interpreter_environment);
\end{verbatim}

変数参照コードは \texttt{Fassq} (コンスセルのcarを探す関数)を使うので、裸のシンボルには絶対にヒットしない。一方 \texttt{let} 側は \texttt{Fmemq} (リストの要素そのものが \texttt{eq} で一致するか探す関数)でこの裸シンボルを検出し、レキシカル束縛をやめて動的束縛に切り替える。

\begin{itemize}
\item 実際の束縛エントリ \texttt{(var . val)} はコンスセルなので、 \texttt{var} そのものとは \texttt{eq} にならない → \texttt{Fmemq} には引っかからない
\item 裸のシンボル \texttt{var} は \texttt{var} そのものなので → \texttt{Fmemq} に引っかかる
\end{itemize}

\textbf{同じ1本のリストの中に、意味の違う2種類のエントリを型の違いだけで共存させている} わけで、なるほどLispらしいと唸った。 \texttt{README.org} の中で \texttt{(defvar my-var)} とだけ書いて後から \texttt{setq} する、というのをよく見かけるが、あれはまさにこの仕組みで「レキシカルスコープのファイル内でも、このシンボルだけは意図的に動的変数として扱わせる」というテクニックだったことになる。
\subsubsection{関数呼び出しの中身: apply\_lambdaとfuncall\_lambda}
\label{sec:org327203e}

関数呼び出し(SUBRでもマクロでもない、通常のクロージャやラムダ呼び出し)は \texttt{apply\_lambda} → \texttt{funcall\_lambda} という2段構えになっている。

\texttt{apply\_lambda} は未評価の引数リストを評価してCの配列に詰めるだけの前処理で、本体は \texttt{funcall\_lambda} にある。ここでは関数の型によって処理が4通りに分かれる。

\begin{itemize}
\item 生の \texttt{(lambda (args) body)} 形式のコンスセル → 動的束縛前提
\item バイトコンパイル済みのレキシカルクロージャ → 引数束縛すら経ずに \texttt{exec\_byte\_code} へ直行
\item インタプリタ実行のクロージャ → コンスタントスロットからレキシカル環境を取得
\item ネイティブコンパイル関数 → \texttt{lambda\_list} から引数名を取得
\end{itemize}

ここでも「レキシカルバイトコードだけは引数束縛ループを一切通らずバイトコードエンジンに丸投げする」という最適化があって、バイトコード自身がスタックから引数を取る命令を持っているから、C側での束縛処理がそもそも不要という理屈だ。

そのあとの \texttt{\&optional} / \texttt{\&rest} を処理する引数束縛ループも中身は素朴で、仮引数リストを1つずつ辿りながら \texttt{\&rest} なら残り全部をリスト化、 \texttt{\&optional} 中で引数が足りなければ \texttt{nil} を補う、という具合。そして最後に束縛するとき、 \texttt{let} のときとまったく同じロジック(\texttt{lexenv} が非nilならconsしてレキシカル、そうでなければ \texttt{specbind} で動的)が繰り返される。つまり「引数の束縛」と「 \texttt{let} の束縛」は、Emacsの内部では同じ1つの判断ロジックのコピーなのだ。
\subsubsection{specpdlという名の何でも屋スタック}
\label{sec:orgb13633f}

動的束縛(\texttt{specbind})がどう元に戻るのかを追うと、 \texttt{specpdl} という単一のスタックに行き着く。 \texttt{lisp.h} の \texttt{specbind\_tag} 列挙型を見ると、このスタックが担っているのは変数の巻き戻しだけではないことが分かる。

\begin{itemize}
\item 動的束縛(通常のletバインディングと、バッファローカル変数用のバリエーション)
\item \texttt{unwind-protect} のクリーンアップ関数
\item バックトレース情報
\item モジュール(動的ロード拡張)のランタイムや環境
\end{itemize}

性質の異なるものを1本のスタックに統一し、 \texttt{unbind\_to} を1回呼ぶだけで「変数を元に戻す」「クリーンアップ関数を呼ぶ」「バックトレースを巻き戻す」が同時に処理される。C言語には例外機構がないので、 \texttt{condition-case} や \texttt{unwind-protect} の巻き戻しは \texttt{setjmp} / \texttt{longjmp} とこの単一スタックの組み合わせだけで実現されている。
\subsubsection{autoloadは「関数の中身」に仕込まれた偽物}
\label{sec:orgc556929}

意外だったのが \texttt{autoload} の実装だ。シンボルの関数セルに \texttt{(autoload FILE ...)} というプレースホルダーが入っていると、 \texttt{eval\_sub} はファイルをロードしてから同じ評価を \texttt{goto retry} でやり直す。

関数を呼ぶたびに「これは本物の関数か、それとも遅延ロードの目印か」を毎回チェックする設計は、標準ライブラリを起動時に全部ロードしないための仕組みで、Emacsの起動速度を支える地味な根幹になっている。 \texttt{straight.el} 経由で入れているパッケージの大半は、初回呼び出しまでこの \texttt{autoload} プレースホルダーのまま待機している。
\subsubsection{40年分の後方互換性という重み}
\label{sec:org1da72fd}

こうして読んでいくと、Emacsが「正しい言語設計」より「既存コードを壊さないこと」を一貫して優先してきたのがよく分かる。1985年当時のダイナミックスコープ前提のコードを今も動かしながら、レキシカルスコープもネイティブコンパイルもバイトコードエンジンも生やしていく。1つの変数が「シンボル単位で動的かレキシカルかを選べる」というのも、1つのalistが「型によって2つの意味を持つ」というのも、本来なら言語設計としては褒められたものではないかもしれない。しかし、後方互換性という制約の中でここまで機能を積み増してきた継ぎ足しの跡が、そのまま設計として残っているのを見るのは素直に楽しい。

次はバイトコードエンジン(\texttt{bytecode.c})か、 \texttt{condition-case} / \texttt{catch}-\texttt{throw} の非局所脱出まわりを読んでみようと思っている。
\subsection{{\bfseries\sffamily DONE} keyboxdのSQLiteエラーからTouch ID復号まで――GPGまわりを一日かけて掘った話}
\label{sec:org5afcb48}
\subsubsection{ある朝、.authinfo.gpg が開けなくなった}
\label{sec:org23f73c1}

きっかけは単純だった。Emacsでいつも通り \texttt{\textasciitilde{}/.authinfo.gpg} を開こうとしたら、こんなエラーが出た。

\begin{verbatim}
Error while decrypting with "/opt/homebrew/bin/gpg":
gpg: keydb_search failed: I/O error during SQL operation
gpg: encrypted with ECDH key, ID D8E725A8316E4066
gpg: key "AA2C51685BA2F0416C2FEB16BEBE9731AC8B2551" not found: I/O error during SQL operation
gpg: all values passed to '--default-key' ignored
gpg: keydb_search failed: I/O error during SQL operation
gpg: public key decryption failed: No secret key
gpg: decryption failed: No secret key
\end{verbatim}

"I/O error during SQL operation" ってなんだ。GPGでSQLのエラーが出るのを初めて見た。
\subsubsection{keyboxdとSQLite、そしてGnuPG 2.5系という沼}
\label{sec:org2ee394e}

調べていくと、GnuPGは2.4系くらいから公開鍵の管理を \texttt{keyboxd} というデーモン経由でSQLiteバックエンド(\texttt{pubring.db})に持たせる方式に移行しつつある。自分の環境は \texttt{brew list -{}-{}versions gnupg} で確認したところ \texttt{2.5.21} 。これはdevel版ではなく、Homebrew-coreが配布している現行の安定版らしい（GnuPGは2025年にバージョニング体系を変えていて、2.5.x系列がメインラインになっている）。

一つずつ潰していった。

\begin{itemize}
\item ロックを握っているプロセスがいないか（ \texttt{lsof} → 何も出ない）
\item \texttt{pubring.db} のパーミッションは \texttt{600} で問題なし
\item \texttt{xattr -l} で見ると \texttt{com.apple.provenance} だけ。怪しいquarantine属性ではない
\item \texttt{sqlite3 pubring.db "PRAGMA integrity\_check;"} → \texttt{ok}
\item ボリュームはローカルAPFS。ネットワークドライブでも同期ストレージでもない
\end{itemize}

ファイルは壊れていない。ロックもない。パーミッションも正常。それでも "I/O error during SQL operation" が出続ける。

こういうとき一番早いのは、原因を完全に特定することより「迂回する」ことだったりする。
\subsubsection{pubring.kbxへの切り戻しで即解決}
\label{sec:orga03ff06}

keyboxdのSQLiteパスをそもそも通らないようにする。レガシー形式の \texttt{pubring.kbx} に戻すだけだ。

\begin{verbatim}
gpgconf --kill all
sed -i '' '/^use-keyboxd$/d' ~/.gnupg/common.conf
mv ~/.gnupg/public-keys.d ~/.gnupg/public-keys.d.broken.$(date +%Y%m%d)
gpgconf --kill all
gpg --list-keys
\end{verbatim}

\begin{verbatim}
gpg: keybox '/Users/ac1965/.gnupg/pubring.kbx' が作成されました
\end{verbatim}

これで "I/O error during SQL operation" は綺麗さっぱり消えた。もちろん公開鍵は空っぽになるので、その後は秘密鍵（ \texttt{private-keys-v1.d} は無傷で残っている）に対応する公開鍵を再インポートする作業が必要になる。

まだ枯れきっていないバックエンドを無理に使い続けるより、確実に動く方に倒す。今回はそれが正解だった。
\subsubsection{ついでにTouch ID化した}
\label{sec:org167c128}

.authinfo.gpg を開くたびにパスフレーズを打つのも面倒だったので、この機会に \texttt{pinentry-touchid} を導入してTouch ID認証に置き換えた。 \texttt{gpg-agent.conf} はこうなっている。

\begin{verbatim}
allow-loopback-pinentry
default-cache-ttl 1800
# allow-emacs-pinentry
pinentry-program /opt/homebrew/bin/pinentry-touchid
\end{verbatim}

\texttt{allow-emacs-pinentry} はあえてコメントアウトしてある。有効にするとEmacsのミニバッファ入力が優先されるケースがあり、Touch ID方式と競合するからだ。
\begin{enumerate}
\item 共通鍵暗号のファイルをTouch IDで復号するときの注意
\label{sec:org3da6a7f}

ここでちょっと引っかかった。 \texttt{pinentry-touchid} は「公開鍵の秘密鍵」と「共通鍵（パスフレーズだけの \texttt{gpg -c}  ）暗号化ファイル」を区別しない。仕組みを見ると、gpg-agentが渡してくる cache-id をキーにmacOSのKeychain（サービス名 \texttt{GnuPG} ）を検索しているだけで、キー種別には興味がない。

つまり共通鍵暗号のファイルにTouch IDを使うと、 \textbf{そのファイルを復号するためのパスフレーズそのものがKeychainに保存される} ことになる。Touch IDは「読み出しのゲート」であって、パスフレーズがSecure Enclaveに守られるわけではない。公式READMEにもそう明記されていた。

はじめは「Touch ID＝安全」と直感的に思ってしまいそうになるが、それは半分正しくて半分違う。以下、実際に運用する上で気をつけていること。

\begin{itemize}
\item \texttt{default-cache-ttl} はgpg-agent自身のメモリキャッシュにしか効かない。Keychainのエントリは手動削除しない限り無期限に残る
\item \texttt{gpg -{}-{}symmetric} は暗号化のたびにS2K saltが変わるので、同じパスフレーズでもファイルが変わるとcache-idも変わり、Keychain未登録扱いになって再度手動入力が必要になることがある。だから毎回中身が変わるファイルより、頻繁に開く1つのファイル（ \texttt{.authinfo.gpg} みたいな）に使うのが現実的
\item どうしても気になる場合は、共通鍵暗号ではなく公開鍵暗号に切り替えるのが筋がいい。その場合Touch IDがガードするのは「秘密鍵へのアクセス」であって、秘密鍵自体は \texttt{\textasciitilde{}/.gnupg/private-keys-v1.d} にファイルとして存在する。パスフレーズをKeychainに丸ごと預けるのとは性質が違う
\end{itemize}

結論として、今回は復号頻度の高い \texttt{.authinfo.gpg} 用途だと割り切ってTouch ID化した。重要度の高いファイルは引き続き公開鍵方式のままにしてある。
\item Keychainへの登録、実は pinentry-mac の仕事
\label{sec:orgf593712}

ここでちょっと訂正しておきたいことがある。「Touch IDでパスフレーズをKeychainに保存する」と書いたが、正確には \texttt{pinentry-touchid} 自身はKeychainへの書き込みを一切やっていない。書き込み・登録の実体は \texttt{pinentry-mac} が持っていて、 \texttt{pinentry-touchid} はそこに登録済みのエントリをTouch ID認証つきで読みに行くだけ、という役割分担になっている（公式READMEに明記されている）。

なので実際にKeychainへ登録する一連の流れはこうなる。

\begin{verbatim}
# 事前に、pinentry-macにKeychain保存を許可しておく
defaults write org.gpgtools.common UseKeychain -bool yes
\end{verbatim}

その状態でパスフレーズ入力が必要な操作（ \texttt{gpg -as -} でも実際の復号でもいい）を行うと、 \texttt{pinentry-mac} のダイアログに「Save in Keychain」のチェックボックスが出る。ここにチェックを入れると、続けてmacOSのログインパスワード入力を求められることがある。これは \texttt{pinentry-mac} や \texttt{pinentry-touchid} がこのエントリにアクセスする権限をOSから正式に認可してもらうためのもので、ここで \textbf{Always Allow} を選んでおく。

さらに \texttt{pinentry-touchid} 側でも初回のTouch ID認証時に、同じように「このエントリへのアクセスを許可するか」を聞かれることがある。README曰く：

\begin{quote}
If a password entry is found but is not "owned" by the pinentry-touchid program after the successful authentication with Touch ID, a normal password will be shown. [\ldots{}] In this dialog click Always allow after entering the password.
\end{quote}

ここでも Always Allow。二段階の認可を経て、ようやく完全にTouch IDだけで完結するようになる。

登録できたかどうかはCLIから確認できる。

\begin{verbatim}
security find-generic-password -s 'GnuPG'
\end{verbatim}

属性の一覧が返れば登録済み。何も見つからなければ次のようなエラーになる。

\begin{verbatim}
security: SecKeychainSearchCopyNext: The specified item could not be found in the keychain.
\end{verbatim}

パスフレーズの権限が完全に \texttt{pinentry-touchid} 側に移った後は、 \texttt{pinentry-mac} 側からの新規書き込みを止めておくとよい。README側も推奨している設定。

\begin{verbatim}
defaults write org.gpgtools.common DisableKeychain -bool yes
\end{verbatim}

エントリを消したいとき（パスフレーズをローテートした時など）は、これも \texttt{security} コマンド一発。

\begin{verbatim}
security delete-generic-password -s 'GnuPG' -a '<アカウント名>'
\end{verbatim}

\texttt{<アカウント名>} は \texttt{find-generic-password -s 'GnuPG'} で返ってくる属性一覧の中の \texttt{acct} の値。gpg-agentが渡すcache-idがベースになっているので、事前に確認してから使う。

「Touch IDが魔法のように守ってくれている」わけではなく、中身は地道な \texttt{pinentry-mac} ↔ Keychain ↔ \texttt{pinentry-touchid} の権限委譲の積み重ねだと分かると、何かあったときの切り分けもしやすくなる。



\begin{itemize}
\item \texttt{gpg-agent} まわりの "I/O error during SQL operation" は、keyboxd/SQLiteバックエンド固有の問題である可能性が高い。ロック・パーミッション・xattr・ファイルシステムをひと通り疑ってダメなら、素直に \texttt{pubring.kbx} に切り戻すのが早い
\item \texttt{pinentry-touchid} は便利だが、内部的には「Touch IDでKeychainのエントリを読み出しているだけ」という理解を持って使う。特に共通鍵暗号ファイルに使う場合は、パスフレーズ自体がKeychainに常駐する点を踏まえて、対象ファイルの重要度を見極める
\end{itemize}

トラブルシューティングは大体こういう順番で潰していくと迷わない。

\begin{enumerate}
\item ロック（他プロセスが握っていないか）
\item パーミッション・所有者
\item ファイル自体の整合性（今回なら \texttt{PRAGMA integrity\_check} ）
\item 環境要因（xattr、ファイルシステム、同期ストレージ）
\item それでもダメならバックエンド自体を疑う
\end{enumerate}

同じ沼にハマった人の参考になれば。
\end{enumerate}
\subsection{{\bfseries\sffamily DONE} 深追いしたら3ヶ月前の全記事画像消失にたどり着いた話}
\label{sec:orgfa41ef5}
きっかけは本当に些細なことだった。ブログのカバー画像が、なんだかいつも
同じ写真になっている気がする——最初はそんな、気のせいかもしれない程度の
違和感だった。Claudeに「ちょっとファクトチェックして」と軽く頼んだ
つもりが、気づけば3ヶ月半前に全記事の画像が消えていたという事実に
たどり着くことになった。以下、その顛末を記録として残しておく。
\subsubsection{ハッシュ値が同じ、という不自然さ}
\label{sec:org0e27634}

Claudeに調査を手伝ってもらいながら \texttt{resources/\_gen} （Hugoのビルドキャッシュ）
を覗いていて、妙なことに気づいた。

\begin{verbatim}
resources/_gen/images/post/2010-1620-4255/cover_hu_be95bd4d5cc45593.jpg
resources/_gen/images/post/2010-3be9-c756/cover_hu_be95bd4d5cc45593.jpg
resources/_gen/images/post/2016-d5bf-eafd/cover_hu_be95bd4d5cc45593.jpg
resources/_gen/images/post/2023-76f0-5099/cover_hu_be95bd4d5cc45593.jpg
\end{verbatim}

Hugoのキャッシュファイル名についているハッシュ値は、元画像のバイト列と
処理オプションから決まる。つまり \textbf{公開日がバラバラな記事間でこのハッシュが
完全一致している} ということは、それぞれの記事に置かれていた
\texttt{cover.jpg} が、物理的にまったく同じファイルだったということになる。

2010年の記事と2023年の記事で、同じ画像。これはさすがに偶然ではない。

このタイミングで、以前から疑っていた \texttt{deploy.sh} の \texttt{post\_month()} 関数の
バグが確信に変わった。

\begin{verbatim}
date_line="$(grep -m1 -E '^(date|publishDate)[[:space:]]*[:=]' "${md_file}" 2>/dev/null || true)"
...
if [[ -z "${month:-}" ]]; then
    month="$(date +%m)"   # ← ここ
    warn "front matterにdateが見つからないため実行時月(${month})にフォールバック"
fi
\end{verbatim}

front matterの日付を取得できなかった場合、 \textbf{実行時点の月} にフォールバック
するようになっていた。1回のデプロイ実行はほぼ同じ時刻に行われるので、
grepが失敗する記事は全部「その日の月」に丸められる。結果、公開日に関係なく
同じ季節画像がコピーされ続けていた、というわけだ。

原因は分かった。あとは直すだけ——のはずだった。
\subsubsection{「復元しよう」と思ったら、そもそも消えていた}
\label{sec:org31fc94f}

カスタムで設定していたはずのカバー画像がいくつか見当たらないことに
気づいて、Gitの履歴を辿ってみることにした。

\begin{verbatim}
for f in content/post/*/cover.jpg; do
  slug=$(dirname "$f" | xargs basename)
  n=$(git log --follow --format=%H -- "$f" | while read c; do
        git show "${c}:${f}" 2>/dev/null | git hash-object --stdin
      done | sort -u | wc -l)
  echo "${slug}: ${n} unique blob(s)"
done
\end{verbatim}

ほとんどの記事で blob 数が2〜4。つまり中身が何度か変わっている。
ここまでは季節ローテーションのバグと辻褄が合う。ところが、全記事の
最新履歴だけを見ると、ある1つのコミットハッシュに揃っていることに
気づいた。

\begin{verbatim}
git log --since="2026-03-20" --until="2026-03-22" \
  --format='%H %ad %s' --date=short -- 'content/post/*/cover.jpg' \
  | sort -u -k1,1
\end{verbatim}

\begin{verbatim}
dea1ff4b49445799379f6183ff758593d924cb82 2026-03-21 Update
\end{verbatim}

このコミット、何をしたのか見てみた。

\begin{verbatim}
git show --stat dea1ff4b49445799379f6183ff758593d924cb82 | head -60
\end{verbatim}

出てきたのは、こういう行がひたすら続く出力だった。

\begin{verbatim}
content/post/2010-1620-4255.md              |  18 -
content/post/2010-1620-4255/cover.jpg       | Bin 1820880 -> 0 bytes
content/post/2010-4ca5-ca76.md              | 744 ---------------------
content/post/2010-4ca5-ca76/cover.jpg       | Bin 65943 -> 0 bytes
...
\end{verbatim}

\texttt{Bin XXXXX -> 0 bytes} 。 \texttt{delete mode} 。全部「削除」だった。

3月21日、46記事のうち、ほぼ全部の \texttt{.md} ファイルと画像ファイル
（カバー画像だけでなく、記事本文中に貼っていた写真も含めて）が、
たった1回のコミットでまとめて消えていた。

「カバー画像のローテーションがおかしい」と思って調べ始めたら、
その裏でもっと大きな損失が3ヶ月半も気づかれずに転がっていた——
というのが、この日いちばん血の気が引いた瞬間だった。
\subsubsection{しめに}
\label{sec:org98d156a}

3月から仕事が忙しくて、ブログのメンテナンスがほとんどできていなかった。
その間に静かに進行していた事故に、こんな回り道をして初めて気づく
ことになるとは思っていなかった。放置していた期間の長さがそのまま
被害の大きさに直結するのだな、と実感した出来事でもあった。

すっきり直って、正直ほっとしている。次はもう少しこまめに見ておこうと思う。
\subsection{{\bfseries\sffamily DONE} カバー画像を直したら、また別の穴が3つ出てきた話}
\label{sec:org4d6f50d}
前回「3ヶ月前の全記事画像消失にたどり着いた話」を書いたとき、最後に「すっきり直って、
正直ほっとしている」と締めた。あれは嘘ではなかったのだけど、正確には「季節ローテーション
のバグだけは直った」であって、話はまだ終わっていなかった。今回はその続き。
\subsubsection{まずは機械的に仕分けるところから}
\label{sec:orgd572e13}

季節画像なのか、記事固有のカバー（以下customと呼ぶ）なのかを目視で1件ずつ確認するのは
さすがに現実的じゃないので、Hugoのビルドキャッシュ \texttt{resources/\_gen/images/post} に残る
ハッシュ値を手がかりに機械的に分類することにした。考え方はシンプルで、

\begin{itemize}
\item 同じハッシュが複数記事に跨って出現する ＝ 使い回されている季節/デフォルト画像
\item 1記事にしか出現しないハッシュ ＝ その記事固有のcustom画像
\end{itemize}

これで46記事を仕分けたところ、custom判定が11件、残り35件が季節側という結果になった。
実際に確認してみると、季節側はさらに2つのグループに分かれていて、片方（9件）は公開月が
ちゃんと同じ季節グループに丸められた、正常なローテーションだった。もう片方（26件）は
公開月がバラバラなのに全部同じ画像を指していて、しかもその画像は今の \texttt{seasonal/01〜04.jpg}
のどれとも一致しない——つまり、前回の記事で解析したあのバグの残骸そのものだった。

該当する26記事の \texttt{cover.jpg} を一旦削除して、修正済みの \texttt{deploy.sh} に再割当てさせたら、
月ごとにきれいに4グループへ分散した。ここまでは想定通り。
\subsubsection{穴その1: 0バイトの \texttt{cover.png} という、また別の残骸}
\label{sec:org7448b52}

custom判定の11記事を \texttt{assets/img/cover/custom/} へ移し終えて、いざ \texttt{hugo server} で
確認したら、カテゴリページのビルドがエラーで落ちた。

\begin{verbatim}
error calling Resize: failed to resize image ".../cover.png": resize ...: image: unknown format
\end{verbatim}

原因は、該当記事のディレクトリに **0バイトの \texttt{cover.png} という残骸**が同居していたこと。
Hugoのページバンドルは \texttt{cover.*} という名前のファイルをfront matterより優先して拾う
慣習があるらしく、以前は同居していた本物の \texttt{cover.jpg} のおかげでたまたま問題が
表面化しなかっただけだった。custom移行で \texttt{cover.jpg} を削除した瞬間、隠れていた
0バイトの残骸が単独で表に出てきた、というオチ。バックアップを取ってから削除して解決。
\subsubsection{穴その2: 分類手法そのものに見落としがあった}
\label{sec:org55fc695}

ところが、これで一件落着かと思いきや、前回の分類に \textbf{入っていない記事} からも同じ
エラーが出た。調べてみると、原因は本物の（0バイトではない）PNG画像。ちゃんと中身の
ある、正真正銘のcustom画像だった。

なぜ最初の分類で漏れたのか。 \texttt{resources/\_gen} はHugoが**実際にレンダリングしたページ
の分しか**キャッシュを作らない。つまり、まだ一度もビルド対象になっていなかった
ページは、キャッシュベースの分類だとそもそも土俵に上がらない。この時点で分類手法を
「Hugoキャッシュ経由」から「 \texttt{content/post/*/cover.*} を拡張子問わず直接ハッシュ
比較」に切り替えて、レンダリング状況に依存しない形に直した。
\subsubsection{穴その3: 自作プロパティ名は、ox-hugoにとってただの飾り}
\label{sec:org2ba4a47}

ここまでの調整で「custom画像をfront matterで明示指定する」仕組み自体は組み上がって
いたのだけど、いざ動かしてみると効かない。 \texttt{all-posts.org} 側には確かに指定を
書き込んだのに、書き出された \texttt{.md} のfront matterに一向にその項目が出てこない。

原因は間抜けなもので、 \texttt{:EXPORT\_HUGO\_COVER:} という \textbf{自分で思いついたプロパティ名} を
そのまま使っていたこと。ox-hugoは \texttt{EXPORT\_FILE\_NAME} や \texttt{EXPORT\_HUGO\_TAGS} のような、
あらかじめ決められた既知のプロパティしかfront matterに変換してくれない。任意の
キーを追加したいなら、 \texttt{EXPORT\_HUGO\_CUSTOM\_FRONT\_MATTER} というplist形式の受け皿を
使うのが正解で、思いつきで名前を作っても黙って無視されるだけだった。

\begin{verbatim}
;; NG: ox-hugoが認識せず、front matterに一切出てこない
:EXPORT_HUGO_COVER: custom/foo.jpg

;; OK: 唯一の任意キー用の経路
:EXPORT_HUGO_CUSTOM_FRONT_MATTER: :cover "custom/foo.jpg"
\end{verbatim}

Emacsのcaptureテンプレート側も同じ間違いをしていたので、新規記事作成時に自動挿入
される空のプレースホルダーも \texttt{:cover ""} の形に直した。
\subsubsection{しめに}
\label{sec:org6e4ee3c}

前回は「3ヶ月半も気づかれずに進行していた事故」がオチだったけど、今回は逆に
「1つ直したら、もぐら叩きみたいに次から次へと別の問題が顔を出す」タイプの話だった。
機械的な分類、ページバンドルの優先順位、ox-hugoの仕様——それぞれは独立した小さな
落とし穴なのだけど、1つの作業の中で3つ続けて踏むと、さすがに笑ってしまった。
\subsection{{\bfseries\sffamily DONE} Blowfish（Hugoテーマ）のショートコード使い方まとめ}
\label{sec:orgb8fda35}
参照元は \href{https://github.com/nunocoracao/blowfish}{nunocoracao/blowfish} の exampleSite ドキュメントと \href{https://blowfish.page/docs/shortcodes/}{公式サイトのshortcodesページ}（2026年7月確認）。

ox-hugoでタングル/エクスポートする前提なので、ショートコードを書く部分は全部 \texttt{\#+begin\_export hugo} 〜 \texttt{\#+end\_export} で囲んである。org地の文（説明・表）はそのままMarkdownとして出力されるので、export ブロックは不要。

org-mode的には「構文そのものを文字として見せたいexportブロック」と「実際にHugoに変換させて機能させるexportブロック」は区別が必要である。 \texttt{\#+begin\_export hugo} はox-hugoに素通しされて実際にレンダリングされてしまうので、構文自体を文字として見せたい場合は \texttt{\#+begin\_example=（または =\#+begin\_src org=）で一段リテラル化してから、その直後に本物の =\#+begin\_export hugo} ブロックを続ける、という2段構成にしている。

ただし \texttt{\#+begin\_example} で囲んだだけでは不十分だった。Hugoのショートコード展開は、囲んでいるのがMarkdownのコードフェンスかどうかを見ずに、ファイル全体のテキストから波括弧2つのショートコード記法を機械的に探して実行してしまうため、コードブロックの中に書いた例文もそのまま実行されてビルドエラーや取得失敗の原因になる。これを防ぐには、Hugo公式のショートコードのエスケープ記法を使う必要がある。開き括弧の直後と閉じ括弧の直前に「/* 」「 */」を挟むと、実行されずに文字通りの文字列としてそのまま出力される。以下すべての節の「書き方はこう」の例では、この記法で統一した。
\subsubsection{Alert（警告・注意喚起ボックス）}
\label{sec:orga7f6a7a}

記事内に目立つメッセージボックスを表示する。中身は Markdown 記法。

書き方はこうなる。

\begin{verbatim}
#+begin_export hugo

**Warning!** This action is destructive!

#+end_export
\end{verbatim}

実際にレンダリングすると、こう表示される。

アイコン付き・色指定もできる。

\begin{verbatim}
#+begin_export hugo

This is an error!

#+end_export
\end{verbatim}

\begin{center}
\begin{tabular}{ll}
パラメータ & 説明\\
\hline
icon & 任意。左側アイコン名。デフォルト: triangle-exclamation\\
iconColor & 任意。アイコン色（hex or カラー名）\\
cardColor & 任意。カード背景色\\
textColor & 任意。テキスト色\\
\end{tabular}
\end{center}
\subsubsection{Admonition（GitHub/Obsidian 風コールアウト）}
\label{sec:org7cf7b74}

これはショートコードではなく Hugo の render hook 経由のMarkdown拡張構文。 \texttt{\#+begin\_export hugo} で囲む必要はなく、そのままorg内でblockquote記法として書ける（ただしox-hugo側でblockquoteへの変換確認が必要。確実に反映したい場合はexportブロック推奨）。

書き方はこう。

\begin{verbatim}
#+begin_export hugo
> [!TIP]
> A Tip type admonition.

> [!TIP]+ Custom Title + Custom Icon
> A collapsible admonition with custom title.
{icon="twitter"}
#+end_export
\end{verbatim}

実際にレンダリングすると、こう表示される。

対応タイプはGitHub系（NOTE / TIP / IMPORTANT / WARNING / CAUTION）とObsidian系（note / abstract / info / todo / tip / success / question / warning / failure / danger / bug / example / quote）。=+=/=-= で折りたたみ制御もできる（Obsidian互換）。
\subsubsection{Accordion（アコーディオン／折りたたみパネル）}
\label{sec:orga2c0171}

書き方はこう。

\begin{verbatim}
#+begin_export hugo

  
  本文（Markdown）
  

  
  インラインアラート
  

#+end_export
\end{verbatim}

実際にレンダリングすると、こう表示される。

\begin{center}
\begin{tabular}{ll}
パラメータ & 説明\\
\hline
mode & collapse（単一展開）/ open（複数展開）。デフォルト collapse\\
separated & true で各項目を独立カード表示\\
title & 必須。ヘッダタイトル\\
open & デフォルトで開いた状態にする\\
icon & タイトル前のアイコン\\
align & left / center / right\\
\end{tabular}
\end{center}
\subsubsection{Ansible Galaxy Card}
\label{sec:org8df0471}

書き方はこう。

\begin{verbatim}
#+begin_export hugo


#+end_export
\end{verbatim}

実際にレンダリングすると、こう表示される。

\texttt{role} または \texttt{collection} のどちらか一方を指定する（ビルド時に \texttt{resources.GetRemote} で取得、ブラウザ側では更新されない）。
\subsubsection{Article（記事埋め込み）}
\label{sec:orga13ad2a}

書き方はこう。

\begin{verbatim}
#+begin_export hugo

#+end_export
\end{verbatim}

実際にレンダリングすると、こう表示される。

\begin{center}
\begin{tabular}{ll}
パラメータ & 説明\\
\hline
link & 必須。埋め込み先の .RelPermalink\\
showSummary & 任意。サマリー表示可否\\
compactSummary & 任意。コンパクト表示\\
\end{tabular}
\end{center}

サブフォルダ運用時はパスを含めて指定する必要がある。
\subsubsection{Badge}
\label{sec:org430bb50}

書き方はこう。

\begin{verbatim}
#+begin_export hugo

New article!

#+end_export
\end{verbatim}

実際にレンダリングすると、こう表示される。
\subsubsection{Button}
\label{sec:org3e095f1}

書き方はこう。 \texttt{pageRef} は存在する内部ページのパスを指定する必要があるので、ここでは記法例として \texttt{docs/getting-started} を挙げているが、実際に使う際はサイト内の実在パスに置き換えること（存在しないパスを \texttt{pageRef} に渡すとビルド時に \texttt{REF\_NOT\_FOUND} エラーになりビルド自体が止まる）。

\begin{verbatim}
#+begin_export hugo

Call to action



Get started

#+end_export
\end{verbatim}

\texttt{href} 版（任意URL/アンカー）はそのままレンダリングできる。

\texttt{pageRef} （内部ページ参照）と \texttt{href} （任意URL）が併用時は \texttt{pageRef} 優先。 \texttt{target} / \texttt{rel} も指定できる。
\subsubsection{Carousel（画像カルーセル）}
\label{sec:org138f9a4}

書き方はこう。 \texttt{captions} の中の \texttt{*formatting*} もMarkdown記法（斜体）である点に注意。

\begin{verbatim}
#+begin_export hugo


{{</* carousel images="gallery/*" captions="{01.jpg:First image with *formatting*,02.jpg:Second image}" */>}}
#+end_export
\end{verbatim}

実際にレンダリングすると、こう表示される。

\begin{center}
\begin{tabular}{ll}
パラメータ & 説明\\
\hline
images & 必須。画像名/URLにマッチする正規表現\\
captions & 任意。key:caption 形式リスト（Markdown可）\\
aspectRatio & 任意。デフォルト 16-9\\
interval & 任意。自動スクロール間隔(ms)。デフォルト 2000\\
\end{tabular}
\end{center}
\subsubsection{Chart（Chart.js グラフ）}
\label{sec:orgcd7eb59}

書き方はこう。

\begin{verbatim}
#+begin_export hugo

type: 'bar',
data: {
  labels: ['Tomato', 'Blueberry', 'Banana'],
  datasets: [{ label: '# of votes', data: [12, 19, 3] }]
}

#+end_export
\end{verbatim}

実際にレンダリングすると、こう表示される。

Chart.jsの設定をそのまま記述する。
\subsubsection{Code Importer（外部コード取り込み）}
\label{sec:orgd484a1b}

書き方はこう。=url= は実在するファイルのraw URLを指定する必要があり、ここでのURLは記法説明用のプレースホルダーなので実レンダリングはしていない（プレースホルダーのままビルドすると \texttt{unable to fetch} の警告が出る）。

\begin{verbatim}
#+begin_export hugo

#+end_export
\end{verbatim}

\begin{center}
\begin{tabular}{ll}
パラメータ & 説明\\
\hline
url & 必須。外部コードファイルURL\\
type & シンタックスハイライト言語\\
startLine / endLine & 任意。取り込み範囲\\
\end{tabular}
\end{center}
\subsubsection{Codeberg / Forgejo / Gitea / GitHub / GitLab Card（外部リポジトリカード）}
\label{sec:orgb02cc16}

書き方はこう。

\begin{verbatim}
#+begin_export hugo





#+end_export
\end{verbatim}

実際にレンダリングすると、こう表示される。

各サービスAPIからスター数・フォーク数等をリアルタイム取得表示する。 \texttt{github} のみ \texttt{showThumbnail} がある。 \texttt{gitlab} は \texttt{baseURL} でセルフホスト対応。
\subsubsection{Email（難読化メールリンク）}
\label{sec:orge8d389e}

書き方はこう。

\begin{verbatim}
#+begin_export hugo

#+end_export
\end{verbatim}

実際にレンダリングすると、こう表示される。
\subsubsection{Figure（画像・図表）}
\label{sec:org6098226}

書き方はこう。

\begin{verbatim}
#+begin_export hugo

#+end_export
\end{verbatim}

実際にレンダリングすると、こう表示される。

標準Markdown画像記法でも自動変換されるので、org中では通常の \texttt{[[file:image.jpg]]} リンクをox-hugoが変換する形でも良い。細かい制御（caption/class/nozoom等）が必要な場合だけexportブロックで明示する。

\begin{center}
\begin{tabular}{ll}
パラメータ & 説明\\
\hline
src & 必須。ページリソース→assets/→static/の順で探索\\
alt & 代替テキスト\\
caption & キャプション(Markdown可)\\
class & <img>への追加CSSクラス\\
figureClass & <figure>への追加CSSクラス（ギャラリー用）\\
href/target & リンク先URL・target\\
nozoom & true でズーム機能無効化\\
default & true で標準Hugo figure挙動に戻す\\
\end{tabular}
\end{center}
\subsubsection{Gallery（画像ギャラリー）}
\label{sec:org9f5a0d1}

書き方はこう。

\begin{verbatim}
#+begin_export hugo

  <img src="gallery/01.jpg" class="grid-w33" />
  <img src="gallery/02.jpg" class="grid-w50 md:grid-w33 xl:grid-w25" />

#+end_export
\end{verbatim}

実際にレンダリングすると、こう表示される。

\texttt{class} \texttt{"grid-wXX"} （10\%〜100\%、5\%刻み。33\%/66\%も可）でカラム幅を指定する。
\subsubsection{Gist（GitHub Gist埋め込み）}
\label{sec:org37da7c9}

書き方はこう。

\begin{verbatim}
#+begin_export hugo


#+end_export
\end{verbatim}

実際にレンダリングすると、こう表示される。

引数は [0] GitHubユーザー名 / [1] Gist ID / [2] 任意、Gist内の特定ファイル名。
\subsubsection{Hugging Face Card}
\label{sec:org0b1b1d5}

書き方はこう。

\begin{verbatim}
#+begin_export hugo


#+end_export
\end{verbatim}

実際にレンダリングすると、こう表示される。

\texttt{model} か \texttt{dataset} のどちらか一方を指定する。
\subsubsection{Icon（SVGアイコン）}
\label{sec:orgc3035d9}

書き方はこう。

\begin{verbatim}
#+begin_export hugo

#+end_export
\end{verbatim}

実際にレンダリングすると、こう表示される。

独自アイコンは \texttt{assets/icons/} にSVG配置し、拡張子なしファイル名で参照できる。
\subsubsection{KaTeX（数式）}
\label{sec:orgea2d807}

書き方はこう。

\begin{verbatim}
#+begin_export hugo

\(f(a,b,c) = (a^2+b^2+c^2)^3\)
#+end_export
\end{verbatim}

実際にレンダリングすると、こう表示される。

記事内に1回だけ配置すればページ内の数式が自動レンダリングされる。org自体の\LaTeX{}記法（==）と併用する場合はexportブロック内でエスケープが崩れないか要確認。
\subsubsection{Keyword / KeywordList（キーワード強調）}
\label{sec:org5100b85}

書き方はこう。

\begin{verbatim}
#+begin_export hugo
 *Super* skill 


 Lorem ipsum dolor. 
 **Important** skill 

#+end_export
\end{verbatim}

実際にレンダリングすると、こう表示される。 \texttt{*Super*} （斜体1個）と \texttt{**Important**} （太字2個）が使い分けられている点に注目。
\subsubsection{Lead（リード文強調）}
\label{sec:org9757613}

書き方はこう。

\begin{verbatim}
#+begin_export hugo

When life gives you lemons, make lemonade.

#+end_export
\end{verbatim}

実際にレンダリングすると、こう表示される。
\subsubsection{List（記事一覧表示）}
\label{sec:org96b88fb}

書き方はこう。

\begin{verbatim}
#+begin_export hugo


#+end_export
\end{verbatim}

実際にレンダリングすると、こう表示される。

\begin{center}
\begin{tabular}{ll}
パラメータ & 説明\\
\hline
limit & 必須。表示件数\\
title & 任意。デフォルト Recent\\
cardView & 任意。カード表示\\
where/value & 任意。記事フィルタ条件（Hugoの where クエリに準拠）\\
\end{tabular}
\end{center}
\subsubsection{LTR/RTL}
\label{sec:org17bc894}

書き方はこう。

\begin{verbatim}
#+begin_export hugo
- This is a markdown list.
{{%/* rtl */%}}
- هذه القائمة باللغة العربية
{{%/* /rtl */%}}
#+end_export
\end{verbatim}

実際にレンダリングすると、こう表示される。
\subsubsection{Markdown Importer（外部Markdown取り込み）}
\label{sec:org6e09746}

書き方はこう。こちらも \texttt{url} は記法説明用のプレースホルダーなので、実在するraw URLに置き換えるまでは実レンダリングしない。

\begin{verbatim}
#+begin_export hugo

#+end_export
\end{verbatim}
\subsubsection{Mermaid（図表・ダイアグラム）}
\label{sec:orgf22a273}

書き方はこう。

\begin{verbatim}
#+begin_export hugo

graph LR;
A[Lemons]-->B[Lemonade];
B-->C[Profit]

#+end_export
\end{verbatim}

実際にレンダリングすると、こう表示される。

org標準の \texttt{\#+begin\_src mermaid} ブロックとは別物である点に注意（Hugo側でレンダリングさせたい場合はexport hugoブロックを使う）。
\subsubsection{Swatches（カラーパレット表示）}
\label{sec:orgf54a007}

書き方はこう。

\begin{verbatim}
#+begin_export hugo

#+end_export
\end{verbatim}

実際にレンダリングすると、こう表示される。
\subsubsection{Tabs（タブ切り替え）}
\label{sec:org416eedf}

書き方はこう。

\begin{verbatim}
#+begin_export hugo

    
    ```javascript
    console.log("Hello");
    ```
    
    
    ```python
    print("Hello")
    ```
    

#+end_export
\end{verbatim}

実際にレンダリングすると、こう表示される。

\begin{center}
\begin{tabular}{ll}
パラメータ & 説明\\
\hline
group & 任意。同一グループ名のタブは連動切替\\
default & 任意。デフォルトアクティブタブのラベル\\
label & 必須（tab側）。タブラベル\\
icon & 任意（tab側）。ラベル前アイコン\\
md & 任意（tab側）。false でネストしたショートコード使用可\\
\end{tabular}
\end{center}
\subsubsection{Timeline（タイムライン）}
\label{sec:orge1d71db}

書き方はこう。

\begin{verbatim}
#+begin_export hugo


本文


#+end_export
\end{verbatim}

実際にレンダリングすると、こう表示される。

\begin{center}
\begin{tabular}{ll}
パラメータ(timelineItem) & 説明\\
\hline
md & Markdownとしてレンダリングするか\\
icon & タイムラインビジュアルに使うアイコン\\
header & 見出し\\
badge & 右上バッジテキスト\\
subheader & サブ見出し\\
\end{tabular}
\end{center}
\subsubsection{TypeIt（タイプライター演出）}
\label{sec:org261866b}

書き方はこう。

\begin{verbatim}
#+begin_export hugo

"Frankly, my dear, I don't give a damn."

#+end_export
\end{verbatim}

実際にレンダリングすると、こう表示される。

\begin{center}
\begin{tabular}{ll}
パラメータ & 説明\\
\hline
tag & 描画するHTMLタグ\\
classList & CSSクラスリスト\\
initialString & 初期表示文字列\\
speed & タイピング速度(ms)\\
lifeLike & 人間らしい不規則タイピング\\
startDelay & 開始までの遅延\\
breakLines & 複数文字列を改行表示するか\\
waitUntilVisible & ビューポート内表示時に開始（デフォルト true）\\
loop & ループ再生\\
\end{tabular}
\end{center}
\subsubsection{Video（動画埋め込み）}
\label{sec:org0ebf5ea}

書き方はこう。=caption= の中の \texttt{**説明文**} もMarkdown記法（太字）である点に注意。

\begin{verbatim}
#+begin_export hugo

#+end_export
\end{verbatim}

実際にレンダリングすると、こう表示される。

\begin{center}
\begin{tabular}{ll}
パラメータ & 説明\\
\hline
src & 必須。ローカル/URL\\
poster & 任意。サムネイル画像\\
caption & 任意。Markdownキャプション\\
autoplay/loop/muted/controls/playsinline & 各種再生制御（真偽値）\\
preload & metadata/none/auto\\
start/end & 再生開始・終了秒数\\
ratio & アスペクト比（16/9等）\\
fit & contain/cover/fill\\
\end{tabular}
\end{center}
\subsubsection{Youtube Lite（YouTube軽量埋め込み）}
\label{sec:org8da2595}

書き方はこう。

\begin{verbatim}
#+begin_export hugo

#+end_export
\end{verbatim}

実際にレンダリングすると、こう表示される。

\begin{center}
\begin{tabular}{ll}
パラメータ & 説明\\
\hline
id & 必須。YouTube動画ID\\
label & 任意。ラベル\\
params & 任意。YouTube Player Parameters（\&区切りで複数指定）\\
\end{tabular}
\end{center}
\subsubsection{補足（org運用時の注意点）}
\label{sec:org15f2605}

\texttt{\#+begin\_export hugo ... \#+end\_export} の中身は ox-hugo によって \textbf{無変換のまま素通し} される。org内の \texttt{*bold*} や \texttt{/italic/} のようなorg記法は効かないので、ショートコード内で強調したい場合はMarkdownの \texttt{**bold**} / \texttt{*italic*} をそのまま書く（この記事もその方針で統一している）。

\texttt{all-posts.org} のようにファイル内に複数記事が \texttt{* 見出し} で区切られている場合、ox-hugoの \texttt{\#+hugo\_base\_dir} や各見出しの \texttt{:EXPORT\_FILE\_NAME:} プロパティ設定が別途必要になる。

既存記事に貼り付ける際は、export blockの直前直後に空行を入れないとorg側のパース崩れが起きるケースがあるので注意する。

Admonitionだけはショートコードではなく \texttt{render-blockquote.html} によるMarkdown render hook実装である。

「構文そのものを文字として見せたい」場合は、=\#+begin\_export hugo= をそのまま書くのではなく \texttt{\#+begin\_example} （または \texttt{\#+begin\_src org} ）で一段リテラル化してから、その直後に本物の \texttt{\#+begin\_export hugo} ブロックを続けて実際のレンダリング結果も見せる、という2段構成にしておくと分かりやすい。

ただし \texttt{\#+begin\_example} で囲むだけでは足りない。Hugoのショートコード展開はMarkdownのコードフェンスの中かどうかを見ずに、ファイル全体から波括弧2つのショートコード記法を探して実行してしまうため、例文として書いたつもりの \texttt{pageRef} や外部URL指定がそのまま実行されて \texttt{REF\_NOT\_FOUND} やfetch失敗の原因になることがある。=\#+begin\_example= の中に書くショートコード例は、開き括弧の直後と閉じ括弧の直前に「/* 」「 */」を挟むエスケープ記法（Hugo公式）で書いておく必要がある。
\subsection{{\bfseries\sffamily DONE} 画像消失を直しに行ったら、直すべき場所が4つあった話}
\label{sec:org9e1a272}
前回、3ヶ月半前から全記事の画像が消えていたことに気づいた話を書いた。
原因は分かったので「あとは直すだけ」のつもりだったのだが、実際に手を
付けてみると、直すべき場所は1箇所どころか4箇所あった。しかも直している
最中に、また別の画像を消すというオチまで付いた。備忘として一通り
記録しておく。
\subsubsection{そもそもの根本原因：flatファイルとPage Bundleの取り違え}
\label{sec:org32dd25d}

うちのHugoサイトは、ox-hugoが記事を \texttt{content/post/<slug>.md} という
flatなファイルとして書き出す構成になっていた。カバー画像は
\texttt{content/post/<slug>/cover.jpg} という、記事ファイルと同名の別
ディレクトリに \texttt{deploy.sh} が機械的に配置する、という運用でずっと
やってきた。

これがそもそもの間違いだった。Hugoの世界では \texttt{<slug>.md} と \texttt{<slug>/}
は同名でも無関係な存在で、後者はPage Resourceとして認識されない。
Blowfishテーマのカバー画像自動検出（ \texttt{.Resources.ByType "image"} ）は
Page Bundle（ \texttt{<slug>/index.md} の形）でしか機能しないので、うちの
構成では最初から蚊帳の外だった。 \texttt{deploy.sh} が季節画像をローテー
ションで機械的に置き続けていたのは、いわばこの構造的な欠陥を隠す
ための対症療法だったわけだ。
\subsubsection{直した場所その1：all-posts.org を一括変換}
\label{sec:orgf2f460d}

ox-hugoには \texttt{:EXPORT\_HUGO\_BUNDLE:} という、まさにこのためのプロパティ
がある。全49記事の \texttt{:EXPORT\_FILE\_NAME: <slug>} を

\begin{verbatim}
:EXPORT_HUGO_BUNDLE: <slug>
:EXPORT_FILE_NAME: index
\end{verbatim}

に機械的に変換した。Hugoは \texttt{<slug>.md} でも \texttt{<slug>/index.md} でも
同じ \texttt{/post/<slug>/} に解決するので、パーマリンクへの影響はゼロ。
これは気楽に一括変換していい変更だった。
\subsubsection{直した場所その2： \texttt{deploy.sh} のカバー画像配置ロジック}
\label{sec:org6062ea2}

ここで最初の見落としに気づいた。カバー画像を置く \texttt{place\_covers()}
関数は \texttt{content/post/*.md} を直接グロブしていたので、記事が
\texttt{<slug>/index.md} に変わった瞬間、対象ファイルを1件も見つけられなく
なって沈黙する。グロブを \texttt{*/index.md} に変え、post\_dirの導出も
文字列結合からindex.mdの親ディレクトリ参照に直した。
\subsubsection{直した場所その3と4：新規記事が古い形式に戻ってしまう2つの入り口}
\label{sec:orgb5254a4}

ここまでで「よし直った」と思っていたのだが、もう2箇所あることに
気づいた。新規記事は \texttt{C-c c b} のorg-captureテンプレートから作る
のだが、Emacs設定（ \texttt{orgx-core.el} のベース版と、それを上書きする
\texttt{personal/ac1965.el} の2箇所）の両方とも \texttt{:EXPORT\_FILE\_NAME:
\%(org-id-new)} のままだった。これでは次に書く記事から、また
flat構造に逆戻りしてしまう。

さらにもう1つ。 \texttt{all-posts.org} の冒頭には、ネタをAIに投げて記事の
下書きを生成させるためのプロンプトテンプレートを仕込んであるの
だが、その中のPROPERTIESドロワーの見本も \texttt{EXPORT\_FILE\_NAME: <slug>}
のままだった。ソースコードだけ直して、その生成元になるテンプレート
文章を見落とすところだった。地味だが、これを見落とすと「なぜか
たまに古い形式の記事が紛れ込む」という原因不明のバグを将来の自分に
仕込むところだった。
\subsubsection{そして、また画像を消した}
\label{sec:org3dec7d2}

一通り直したので、動作確認のつもりで一度 \texttt{content/post} を
\texttt{rm -rf} してから再exportしてみた。 \texttt{content/post} はビルド成果物
なので.gitignore対象、消しても再生成されるだけのはずだった。

ところが、記事本文で使っていたBlowfishのfigure/carouselショートコード
（ \texttt{src}"frieren.jpg"= や \texttt{images}"inuyama-fest/*"= のように画像を
直接指定するタイプ）が参照していた画像だけは戻ってこなかった。

理由は単純で、これらはPage Resourceとして \texttt{content/post/<slug>/}
直下に \textbf{手で置いていただけ} のファイルだった。カバー画像は
\texttt{deploy.sh} が毎回決定的に再生成するし、Orgの \texttt{[[file:...]]} リンク
はox-hugoが自動でコピーし直してくれる。しかしshortcode内の
\texttt{src}"\ldots{}"= はox-hugoにとってただの文字列で、どの自動生成の仕組みにも
紐づいていなかった。一番安全なはずの「ビルド成果物だから消していい」
という前提が、実はこの種の画像には最初から成り立っていなかった、
というわけだ。
\subsubsection{deployブランチという名の保険}
\label{sec:org4fead5c}

さすがに焦ったが、幸い逃げ道が残っていた。うちの \texttt{deploy.sh} は
ビルド後の \texttt{public/} を、sourceとは別の \texttt{deploy} ブランチにpushして
いる。つまり、公開されているHTML・画像そのものが独立したGit履歴を
持っているということだ。

GitHubのTrees APIで \texttt{deploy} ブランチの履歴を遡り、3月21日の大量
削除が起きる直前、最後に正常だったdeployコミットを特定した。そこから
Hugoのリサイズキャッシュ（ \texttt{*\_hu\_<hash>.*} ）を除いた \textbf{オリジナル
ファイルだけ} を掘り出したところ、10記事・37ファイル、 \texttt{frieren.jpg}
も \texttt{inuyama-fest/} の写真11枚も無事に発見できた。全ファイル、コミット
時点のGit blobサイズと突き合わせて一致を確認済みで、EXIF情報から
撮影日時まで読み取れる、正真正銘のオリジナルだった。

これを機に、 \texttt{assets/img/post-assets/<slug>/} という「正本置き場」を
新設し、 \texttt{deploy.sh} に \texttt{sync\_post\_assets()} を追加した。build毎に
ここから \texttt{content/post/<slug>/} へ同期するので、以後は \texttt{rm -rf
content/post} してもこの種の画像はもう消えない。
\subsubsection{しめに}
\label{sec:orga9d6c28}

「カバー画像のローテーションがおかしい」から始まった話が、Page
Bundleへの構造変更、Emacs側の2つの入り口、AIプロンプトテンプレート、
そして本文画像そのものの消失発覚と、芋づる式に広がっていった。

一番の教訓は、「ビルド成果物だから消していい」という前提は、何に
よってどう再生成されるかを1つずつ確認しない限り信用してはいけない、
ということだった。cover.jpgとインライン画像は見た目には
同じ「content/post配下のファイル」だが、裏の仕組みはまったく別物
だった。今回はたまたまdeployブランチという保険が効いたからよかった
ものの、無ければ本当に取り返しがつかなかった。

直すべき場所を1つ見つけるたびに、もう1つ見つかる。地味に骨の折れる
作業だったが、これで少なくとも同じ事故は起きなくなったはずだ。
\subsubsection{参考：Page Bundleとは}
\label{sec:orgf457b07}

なぜflatファイルとPage Bundleを取り違えるとこんなことになるのか、
ここで改めて整理しておく。
\subsection{{\bfseries\sffamily DONE} Makefile が自分自身を tangle する設定で、実機でしかハマらないバグを4連発踏んだ話}
\label{sec:org723a813}
前回、 \texttt{README.org} の Makefile ブロックでタブが壊れる話を書いたけど、あの後も
続きがあった。実際に手元で \texttt{make reload} や \texttt{make diagrams} を叩くたびに、
サンドボックス環境の検証だけでは踏めない類のバグが次々出てきて、結局4段階の
実機デバッグになった。せっかくなので記録として残しておく。

\textbf{最初の直しが別のバグを生んだ}

タブが \texttt{org-babel-tangle} に飲み込まれる問題、最初は \texttt{.RECIPEPREFIX :} >= で
recipe行の先頭マーカーをタブ以外の文字に変える案で直したつもりだった。だがこれは
GNU Make 3.82以降の機能で、macOS標準の \texttt{/usr/bin/make} はGPLv3回避のため3.81に
凍結されている。案の定、手元で試してもらったら:

\begin{verbatim}
Makefile:87: *** missing separator.  Stop.
\end{verbatim}

同じエラーが別の行で再発した。バージョン依存の機能に頼ったのが失敗で、結局
本物のタブ文字に戻し、 \texttt{org-src-preserve-indentation} を \texttt{t} にして tangle時の
インデント除去処理自体を無効化する方式に変更した。

\textbf{その修正がまた別のファイルを壊した}

ところがこの \texttt{org-src-preserve-indentation} を \texttt{(setq ...)} で \textbf{グローバルに}
設定したのが次の落とし穴だった。README.org は1回のEmacsセッションで全ブロックを
tangleするので、Makefileだけでなく \texttt{early-init.el} や \texttt{lisp/*.el} 全部に
この設定が波及してしまい、org側の見た目そろえ用スペースが \texttt{early-init.el} の
1行目（ \texttt{lexical-binding} cookie）にまで残ってしまった。結果、Emacs起動時に
こんな警告が出た。

\begin{quote}
Warning (files): Missing 'lexical-binding' cookie in "\textasciitilde{}/.emacs.d-stable/early-init.el".
\end{quote}

直し方は2段階tangleにすること。1段階目は普通に（indentation preserveなし）
全ブロックをtangleして \texttt{early-init.el} 等を正しく生成し、2段階目で
\texttt{org-babel-tangle-file} の言語フィルタ引数に \texttt{\textbackslash{}"makefile\textbackslash{}"} だけを渡して、
Makefileブロックだけを preserve-indentation=t で再tangleする。これで両方
正しくなった。

\textbf{leafマクロが理由不明のまま原因不特定になった件}

Makefile周りが片付いたところで、今度は実際の起動ログから \texttt{ui-which-key}
モジュールが読み込み失敗しているのが見つかった。

\begin{verbatim}
[modules] Failed to load ui-which-key: Symbol's value as variable is void: which-key
\end{verbatim}

該当箇所は \texttt{leaf which-key} を \texttt{:straight} と \texttt{:if} の組み合わせで
Emacs 30以降／未満に条件分岐させていた部分。手元にEmacs 30系が無かったので、
straight・leaf・leaf-keywordsを実際にブートストラップして単体評価まで試したが、
正直なところ \texttt{leaf-keywords.el} の \texttt{:straight} 拡張内部の正確な原因までは
特定しきれなかった。原因究明より実利を取って、 \texttt{leaf} マクロを経由せず
素の \texttt{(when (< emacs-major-version 30) (straight-use-package 'which-key))}
に書き換えることで問題自体を回避した。

\textbf{Puppeteerのバージョン固定にも足をすくわれた}

図版のMermaidレンダリング（ \texttt{mmdc} ）でも似たようなことが起きた。エラーメッセージは
素直に「Chromeが見つからない」と言っていて、Puppeteer公式の案内通り
\texttt{npx puppeteer browsers install chrome-headless-shell} を勧めたのだが、これが
外れだった。 \texttt{mmdc} が同梱する \texttt{puppeteer-core} は特定バージョンのChromeに
ピン留めされているため、システムに最新版を入れても一致せず、同じエラーが
再発し続けた。

\textbf{結局16枚全部揃った}

最終的に \texttt{dot-tangle} / \texttt{dot-svg} / \texttt{mmd-tangle} / \texttt{mmd-svg} を1つにまとめた
\texttt{make diagrams} で、design\_spec.org側13枚 ＋ README.org側3枚の計16枚が
問題なくレンダリングされるところまで確認できた。

サンドボックスでの検証だけでは絶対に踏めないバグ（Makeのバージョン差、
tangleの副作用範囲、leafマクロの内部実装、puppeteerのバージョン固定）が
4段階連続で出てきて、そのたびに「サンドボックスでは動いたのに」という
展開になった。実機で試してもらうまで確信が持てない類の不具合が、こんなに
積み重なるとは思っていなかった
\end{document}
